% latex-rules — reusable preamble (see the latex-rules skill).
% Reusable macros only; define paper-specific notation in your main .tex, not here.

% ==== packages ====

% ---- fonts & encoding ----
\usepackage[T1]{fontenc}   % proper 8-bit fonts + real hyphenation
\usepackage{lmodern}       % scalable Latin Modern (fixes bitmap CM warnings)

% ---- math ----
\usepackage{amsmath}
\usepackage{amssymb}
\usepackage{mathtools}     % loads + extends amsmath
\usepackage{dsfont}        % \mathds{1}

% ---- graphics & drawing ----
\usepackage{graphicx}
\usepackage[export]{adjustbox}
\usepackage{tikz}
\usetikzlibrary{shapes}
\usetikzlibrary{positioning}
\usetikzlibrary{trees}
%\usetikzlibrary{graphs,graphdrawing}
%\usegdlibrary{trees}
\usepackage[most]{tcolorbox}
\usepackage{subcaption}    % NOT subfigure / subfig
\usepackage{pdflscape}

% ---- color ----
\usepackage{xcolor}

% ---- tables ----
\usepackage{booktabs}
\usepackage{multirow}
\usepackage{makecell}
\usepackage{colortbl}

% ---- typography, lists, units ----
\usepackage[protrusion=true,expansion=true]{microtype}  % char protrusion + font expansion; drop expansion under XeLaTeX (unsupported)
\usepackage{csquotes}
\usepackage{enumitem}
\usepackage{siunitx}
\sisetup{detect-all, per-mode=symbol}   % detect-all: match surrounding font; per-mode=symbol: m/s not m s^-1
\usepackage{xspace}

% ---- floats, algorithms, symbols ----
\usepackage{float}
\usepackage{afterpage}
\usepackage{footmisc}
\usepackage{listings}
\usepackage{algorithm}
\usepackage{algpseudocode}   % algorithmicx; modern replacement for algorithmic
\usepackage{pifont}
%\usepackage{pgfplots}

% ---- occasionally useful (uncomment as needed) ----
%\usepackage{tabu}
%\usepackage[defaultlines=3,all]{nowidow}
%\usepackage{setspace}
%\usepackage{geometry}
%\usepackage[all]{xy}
%\usepackage{natbib}
%\usepackage{subfigure}
% \usepackage{svg}   % needs --shell-escape + inkscape and is rejected by arXiv; pre-convert SVGs to PDF instead

% ==== math notation ====
% Symbol macros are wrapped in \ensuremath, so they also work in text (\R, \bx, ... auto-enter math).

% ---- quote marks in math ----
\DeclareMathSymbol{\mlq}{\mathord}{operators}{``}
\DeclareMathSymbol{\mrq}{\mathord}{operators}{`'}

% ---- number sets (blackboard bold) ----
% NOTE: \P overrides the built-in pilcrow (¶); see also \S \O \L \H below.
\newcommand{\R}{\ensuremath{\mathbb{R}}}
\newcommand{\N}{\ensuremath{\mathbb{N}}}
\newcommand{\Z}{\ensuremath{\mathbb{Z}}}
\newcommand{\Q}{\ensuremath{\mathbb{Q}}}
\newcommand{\C}{\ensuremath{\mathbb{C}}}
\newcommand{\D}{\ensuremath{\mathbb{D}}}
\newcommand{\E}{\ensuremath{\mathbb{E}}}
\newcommand{\X}{\ensuremath{\mathbb{X}}}
\newcommand{\Y}{\ensuremath{\mathbb{Y}}}
\renewcommand{\P}{\ensuremath{\mathbb{P}}}
\newcommand{\sigB}{\ensuremath{\mathbb{B}}}
\newcommand{\1}{\ensuremath{\mathds{1}}}

% ---- calligraphic sets ----
% NOTE: \S \O \L \H override built-in text commands (§, Ø, Ł, and the o-double-acute accent \H{o}).
\newcommand{\calN}{\ensuremath{\mathcal{N}}}
\newcommand{\calD}{\ensuremath{\mathcal{D}}}
\newcommand{\A}{\ensuremath{\mathcal{A}}}
\newcommand{\B}{\ensuremath{\mathcal{B}}}
\newcommand{\I}{\ensuremath{\mathcal{I}}}
\newcommand{\F}{\ensuremath{\mathcal{F}}}
\newcommand{\G}{\ensuremath{\mathcal{G}}}
\newcommand{\M}{\ensuremath{\mathcal{M}}}
\newcommand{\Tau}{\ensuremath{\mathcal{T}}}
\newcommand{\W}{\ensuremath{\mathcal{W}}}
\newcommand{\sigE}{\ensuremath{\mathcal{E}}}
\renewcommand{\H}{\ensuremath{\mathcal{H}}}
\renewcommand{\L}{\ensuremath{\mathcal{L}}}
\renewcommand{\S}{\ensuremath{\mathcal{S}}}
\renewcommand{\O}{\ensuremath{\mathcal{O}}}

% ---- named operators ----
% NOTE: \Re \Im are redefined to roman "Re"/"Im" (not the default blackletter ℜ/ℑ).
\DeclareMathOperator{\softmax}{softmax}
\DeclareMathOperator*{\argmax}{arg\,max}   % * = subscript sits underneath in display
\DeclareMathOperator{\Var}{Var}
\newcommand{\cov}{\ensuremath{\operatorname{cov}}}
\newcommand{\corr}{\ensuremath{\operatorname{corr}}}
\newcommand{\spanop}{\ensuremath{\operatorname{span}}}
\newcommand{\att}{\ensuremath{\operatorname{attention}}}
\newcommand{\suff}{\ensuremath{\operatorname{suff}}}
\newcommand{\comp}{\ensuremath{\operatorname{comp}}}
\newcommand{\In}{\ensuremath{\operatorname{In}}}
\newcommand{\opcap}{\ensuremath{\operatorname{cap}}}
\newcommand{\entr}{\ensuremath{\operatorname{entries}}}
\newcommand{\ops}{\ensuremath{\operatorname{ops}}}
\newcommand{\id}{\ensuremath{\mathrm{id}}}
\renewcommand{\Re}{\ensuremath{\operatorname{Re}}}
\renewcommand{\Im}{\ensuremath{\operatorname{Im}}}

% ---- symbols & shorthands ----
\newcommand{\eps}{\ensuremath{\varepsilon}}
\renewcommand{\phi}{\ensuremath{\varphi}}     % deliberate: \phi -> varphi
\newcommand{\del}{\ensuremath{\partial}}
\newcommand{\grad}{\ensuremath{\nabla}}
\newcommand{\indep}{\ensuremath{\perp\!\!\!\perp}}
\newcommand{\emptyword}{\ensuremath{\varnothing}}
\newcommand{\tensor}{\ensuremath{\otimes}}
\newcommand{\bigtensor}{\ensuremath{\bigotimes}}
\newcommand{\To}{\ensuremath{\longrightarrow}}
\newcommand{\eqover}[1]{\ensuremath{\stackrel{#1}{=}}}
\newcommand{\darover}[1]{\ensuremath{\xrightarrow{#1}}}
\newcommand{\Hoel}{\text{Höl}}

% ---- paired delimiters ----
\DeclarePairedDelimiter{\norm}{\lVert}{\rVert}   % \norm{x}, \norm*{x} (full size), \norm[\big]{x}
\DeclarePairedDelimiter{\abs}{\lvert}{\rvert}

% ---- bold vectors / matrices ----
\newcommand{\bx}{\ensuremath{\mathbf{x}}}
\newcommand{\bX}{\ensuremath{\mathbf{X}}}
\newcommand{\by}{\ensuremath{\mathbf{y}}}
\newcommand{\bz}{\ensuremath{\mathbf{z}}}
\newcommand{\bB}{\ensuremath{\mathbf{B}}}
\newcommand{\bp}{\ensuremath{\mathbf{p}}}
\newcommand{\bW}{\ensuremath{\mathbf{W}}}
\newcommand{\bt}{\ensuremath{\mathbf{t}}}
\newcommand{\be}{\ensuremath{\mathbf{e}}}

% ---- small constructs ----
\newcommand{\Ball}[2]{\ensuremath{B_#1(#2)}}   % \Ball{\eps}{center}
\newcommand{\twobytwo}[4]{\ensuremath{\begin{pmatrix} #1 & #2 \\ #3 & #4 \end{pmatrix}}}

% ==== text, color & layout helpers ====
\definecolor{darkgreen}{HTML}{009B55}
\newcommand{\gtxt}[1]{\text{\textcolor{gray}{#1}}}
\newcommand{\grntxt}[1]{\text{\textcolor{darkgreen}{#1}}}
\newcommand{\rdtxt}[1]{\text{\textcolor{red}{#1}}}
\newcommand{\code}[1]{\texttt{#1}}
\newcommand{\tldr}{\textbf{TL;DR:}\xspace}
\newcommand{\parahead}[1]{\noindent\textbf{#1.}\enspace}  % run-in heading; use instead of \paragraph to save space
\newcommand*\rot{\rotatebox{90}}

% ---- table cell helpers ----
\newcommand{\cmark}{\ding{51}}
\newcommand{\xmark}{\ding{55}}
\newcommand{\rowhighlight}{\rowcolor{gray!30}}
\newcommand{\cellhighlight}{\cellcolor{yellow!40}}

% ---- rebuttal reviewer tags: bold, distinct colors to signpost which reviewer a point answers ----
\newcommand{\Rone}{\textbf{\textcolor[HTML]{1F77B4}{R1}}}    % blue
\newcommand{\Rtwo}{\textbf{\textcolor[HTML]{D62728}{R2}}}    % red
\newcommand{\Rthree}{\textbf{\textcolor[HTML]{2CA02C}{R3}}}  % green
\newcommand{\Rfour}{\textbf{\textcolor[HTML]{D9730D}{R4}}}   % muted orange

% ---- takeaway callout box ----
\newtcolorbox{takeaway}[1][Takeaway]{
  colback=blue!4!white,      % very light fill
  colframe=blue!50!black,    % muted dark-blue frame
  boxrule=0.6pt,
  arc=2pt,
  left=6pt, right=6pt, top=2pt, bottom=2pt,
  title=#1,
  fonttitle=\bfseries\small,
  coltitle=white,
  colbacktitle=blue!50!black
}

% ==== links & smart references: load LAST (delete if your class already loads hyperref) ====
\usepackage[hidelinks]{hyperref}
\usepackage[capitalize,noabbrev]{cleveref}   % \cref{...}, \Cref{...} — must come after hyperref

% ==== draft mode & \todo ====
% "Explicit" = the main .tex defined \ifdraftmode before % latex-rules — reusable preamble (see the latex-rules skill).
% Reusable macros only; define paper-specific notation in your main .tex, not here.

% ==== packages ====

% ---- fonts & encoding ----
\usepackage[T1]{fontenc}   % proper 8-bit fonts + real hyphenation
\usepackage{lmodern}       % scalable Latin Modern (fixes bitmap CM warnings)

% ---- math ----
\usepackage{amsmath}
\usepackage{amssymb}
\usepackage{mathtools}     % loads + extends amsmath
\usepackage{dsfont}        % \mathds{1}

% ---- graphics & drawing ----
\usepackage{graphicx}
\usepackage[export]{adjustbox}
\usepackage{tikz}
\usetikzlibrary{shapes}
\usetikzlibrary{positioning}
\usetikzlibrary{trees}
%\usetikzlibrary{graphs,graphdrawing}
%\usegdlibrary{trees}
\usepackage[most]{tcolorbox}
\usepackage{subcaption}    % NOT subfigure / subfig
\usepackage{pdflscape}

% ---- color ----
\usepackage{xcolor}

% ---- tables ----
\usepackage{booktabs}
\usepackage{multirow}
\usepackage{makecell}
\usepackage{colortbl}

% ---- typography, lists, units ----
\usepackage[protrusion=true,expansion=true]{microtype}  % char protrusion + font expansion; drop expansion under XeLaTeX (unsupported)
\usepackage{csquotes}
\usepackage{enumitem}
\usepackage{siunitx}
\sisetup{detect-all, per-mode=symbol}   % detect-all: match surrounding font; per-mode=symbol: m/s not m s^-1
\usepackage{xspace}

% ---- floats, algorithms, symbols ----
\usepackage{float}
\usepackage{afterpage}
\usepackage{footmisc}
\usepackage{listings}
\usepackage{algorithm}
\usepackage{algpseudocode}   % algorithmicx; modern replacement for algorithmic
\usepackage{pifont}
%\usepackage{pgfplots}

% ---- occasionally useful (uncomment as needed) ----
%\usepackage{tabu}
%\usepackage[defaultlines=3,all]{nowidow}
%\usepackage{setspace}
%\usepackage{geometry}
%\usepackage[all]{xy}
%\usepackage{natbib}
%\usepackage{subfigure}
% \usepackage{svg}   % needs --shell-escape + inkscape and is rejected by arXiv; pre-convert SVGs to PDF instead

% ==== math notation ====
% Symbol macros are wrapped in \ensuremath, so they also work in text (\R, \bx, ... auto-enter math).

% ---- quote marks in math ----
\DeclareMathSymbol{\mlq}{\mathord}{operators}{``}
\DeclareMathSymbol{\mrq}{\mathord}{operators}{`'}

% ---- number sets (blackboard bold) ----
% NOTE: \P overrides the built-in pilcrow (¶); see also \S \O \L \H below.
\newcommand{\R}{\ensuremath{\mathbb{R}}}
\newcommand{\N}{\ensuremath{\mathbb{N}}}
\newcommand{\Z}{\ensuremath{\mathbb{Z}}}
\newcommand{\Q}{\ensuremath{\mathbb{Q}}}
\newcommand{\C}{\ensuremath{\mathbb{C}}}
\newcommand{\D}{\ensuremath{\mathbb{D}}}
\newcommand{\E}{\ensuremath{\mathbb{E}}}
\newcommand{\X}{\ensuremath{\mathbb{X}}}
\newcommand{\Y}{\ensuremath{\mathbb{Y}}}
\renewcommand{\P}{\ensuremath{\mathbb{P}}}
\newcommand{\sigB}{\ensuremath{\mathbb{B}}}
\newcommand{\1}{\ensuremath{\mathds{1}}}

% ---- calligraphic sets ----
% NOTE: \S \O \L \H override built-in text commands (§, Ø, Ł, and the o-double-acute accent \H{o}).
\newcommand{\calN}{\ensuremath{\mathcal{N}}}
\newcommand{\calD}{\ensuremath{\mathcal{D}}}
\newcommand{\A}{\ensuremath{\mathcal{A}}}
\newcommand{\B}{\ensuremath{\mathcal{B}}}
\newcommand{\I}{\ensuremath{\mathcal{I}}}
\newcommand{\F}{\ensuremath{\mathcal{F}}}
\newcommand{\G}{\ensuremath{\mathcal{G}}}
\newcommand{\M}{\ensuremath{\mathcal{M}}}
\newcommand{\Tau}{\ensuremath{\mathcal{T}}}
\newcommand{\W}{\ensuremath{\mathcal{W}}}
\newcommand{\sigE}{\ensuremath{\mathcal{E}}}
\renewcommand{\H}{\ensuremath{\mathcal{H}}}
\renewcommand{\L}{\ensuremath{\mathcal{L}}}
\renewcommand{\S}{\ensuremath{\mathcal{S}}}
\renewcommand{\O}{\ensuremath{\mathcal{O}}}

% ---- named operators ----
% NOTE: \Re \Im are redefined to roman "Re"/"Im" (not the default blackletter ℜ/ℑ).
\DeclareMathOperator{\softmax}{softmax}
\DeclareMathOperator*{\argmax}{arg\,max}   % * = subscript sits underneath in display
\DeclareMathOperator{\Var}{Var}
\newcommand{\cov}{\ensuremath{\operatorname{cov}}}
\newcommand{\corr}{\ensuremath{\operatorname{corr}}}
\newcommand{\spanop}{\ensuremath{\operatorname{span}}}
\newcommand{\att}{\ensuremath{\operatorname{attention}}}
\newcommand{\suff}{\ensuremath{\operatorname{suff}}}
\newcommand{\comp}{\ensuremath{\operatorname{comp}}}
\newcommand{\In}{\ensuremath{\operatorname{In}}}
\newcommand{\opcap}{\ensuremath{\operatorname{cap}}}
\newcommand{\entr}{\ensuremath{\operatorname{entries}}}
\newcommand{\ops}{\ensuremath{\operatorname{ops}}}
\newcommand{\id}{\ensuremath{\mathrm{id}}}
\renewcommand{\Re}{\ensuremath{\operatorname{Re}}}
\renewcommand{\Im}{\ensuremath{\operatorname{Im}}}

% ---- symbols & shorthands ----
\newcommand{\eps}{\ensuremath{\varepsilon}}
\renewcommand{\phi}{\ensuremath{\varphi}}     % deliberate: \phi -> varphi
\newcommand{\del}{\ensuremath{\partial}}
\newcommand{\grad}{\ensuremath{\nabla}}
\newcommand{\indep}{\ensuremath{\perp\!\!\!\perp}}
\newcommand{\emptyword}{\ensuremath{\varnothing}}
\newcommand{\tensor}{\ensuremath{\otimes}}
\newcommand{\bigtensor}{\ensuremath{\bigotimes}}
\newcommand{\To}{\ensuremath{\longrightarrow}}
\newcommand{\eqover}[1]{\ensuremath{\stackrel{#1}{=}}}
\newcommand{\darover}[1]{\ensuremath{\xrightarrow{#1}}}
\newcommand{\Hoel}{\text{Höl}}

% ---- paired delimiters ----
\DeclarePairedDelimiter{\norm}{\lVert}{\rVert}   % \norm{x}, \norm*{x} (full size), \norm[\big]{x}
\DeclarePairedDelimiter{\abs}{\lvert}{\rvert}

% ---- bold vectors / matrices ----
\newcommand{\bx}{\ensuremath{\mathbf{x}}}
\newcommand{\bX}{\ensuremath{\mathbf{X}}}
\newcommand{\by}{\ensuremath{\mathbf{y}}}
\newcommand{\bz}{\ensuremath{\mathbf{z}}}
\newcommand{\bB}{\ensuremath{\mathbf{B}}}
\newcommand{\bp}{\ensuremath{\mathbf{p}}}
\newcommand{\bW}{\ensuremath{\mathbf{W}}}
\newcommand{\bt}{\ensuremath{\mathbf{t}}}
\newcommand{\be}{\ensuremath{\mathbf{e}}}

% ---- small constructs ----
\newcommand{\Ball}[2]{\ensuremath{B_#1(#2)}}   % \Ball{\eps}{center}
\newcommand{\twobytwo}[4]{\ensuremath{\begin{pmatrix} #1 & #2 \\ #3 & #4 \end{pmatrix}}}

% ==== text, color & layout helpers ====
\definecolor{darkgreen}{HTML}{009B55}
\newcommand{\gtxt}[1]{\text{\textcolor{gray}{#1}}}
\newcommand{\grntxt}[1]{\text{\textcolor{darkgreen}{#1}}}
\newcommand{\rdtxt}[1]{\text{\textcolor{red}{#1}}}
\newcommand{\code}[1]{\texttt{#1}}
\newcommand{\tldr}{\textbf{TL;DR:}\xspace}
\newcommand{\parahead}[1]{\noindent\textbf{#1.}\enspace}  % run-in heading; use instead of \paragraph to save space
\newcommand*\rot{\rotatebox{90}}

% ---- table cell helpers ----
\newcommand{\cmark}{\ding{51}}
\newcommand{\xmark}{\ding{55}}
\newcommand{\rowhighlight}{\rowcolor{gray!30}}
\newcommand{\cellhighlight}{\cellcolor{yellow!40}}

% ---- rebuttal reviewer tags: bold, distinct colors to signpost which reviewer a point answers ----
\newcommand{\Rone}{\textbf{\textcolor[HTML]{1F77B4}{R1}}}    % blue
\newcommand{\Rtwo}{\textbf{\textcolor[HTML]{D62728}{R2}}}    % red
\newcommand{\Rthree}{\textbf{\textcolor[HTML]{2CA02C}{R3}}}  % green
\newcommand{\Rfour}{\textbf{\textcolor[HTML]{D9730D}{R4}}}   % muted orange

% ---- takeaway callout box ----
\newtcolorbox{takeaway}[1][Takeaway]{
  colback=blue!4!white,      % very light fill
  colframe=blue!50!black,    % muted dark-blue frame
  boxrule=0.6pt,
  arc=2pt,
  left=6pt, right=6pt, top=2pt, bottom=2pt,
  title=#1,
  fonttitle=\bfseries\small,
  coltitle=white,
  colbacktitle=blue!50!black
}

% ==== links & smart references: load LAST (delete if your class already loads hyperref) ====
\usepackage[hidelinks]{hyperref}
\usepackage[capitalize,noabbrev]{cleveref}   % \cref{...}, \Cref{...} — must come after hyperref

% ==== draft mode & \todo ====
% "Explicit" = the main .tex defined \ifdraftmode before % latex-rules — reusable preamble (see the latex-rules skill).
% Reusable macros only; define paper-specific notation in your main .tex, not here.

% ==== packages ====

% ---- fonts & encoding ----
\usepackage[T1]{fontenc}   % proper 8-bit fonts + real hyphenation
\usepackage{lmodern}       % scalable Latin Modern (fixes bitmap CM warnings)

% ---- math ----
\usepackage{amsmath}
\usepackage{amssymb}
\usepackage{mathtools}     % loads + extends amsmath
\usepackage{dsfont}        % \mathds{1}

% ---- graphics & drawing ----
\usepackage{graphicx}
\usepackage[export]{adjustbox}
\usepackage{tikz}
\usetikzlibrary{shapes}
\usetikzlibrary{positioning}
\usetikzlibrary{trees}
%\usetikzlibrary{graphs,graphdrawing}
%\usegdlibrary{trees}
\usepackage[most]{tcolorbox}
\usepackage{subcaption}    % NOT subfigure / subfig
\usepackage{pdflscape}

% ---- color ----
\usepackage{xcolor}

% ---- tables ----
\usepackage{booktabs}
\usepackage{multirow}
\usepackage{makecell}
\usepackage{colortbl}

% ---- typography, lists, units ----
\usepackage[protrusion=true,expansion=true]{microtype}  % char protrusion + font expansion; drop expansion under XeLaTeX (unsupported)
\usepackage{csquotes}
\usepackage{enumitem}
\usepackage{siunitx}
\sisetup{detect-all, per-mode=symbol}   % detect-all: match surrounding font; per-mode=symbol: m/s not m s^-1
\usepackage{xspace}

% ---- floats, algorithms, symbols ----
\usepackage{float}
\usepackage{afterpage}
\usepackage{footmisc}
\usepackage{listings}
\usepackage{algorithm}
\usepackage{algpseudocode}   % algorithmicx; modern replacement for algorithmic
\usepackage{pifont}
%\usepackage{pgfplots}

% ---- occasionally useful (uncomment as needed) ----
%\usepackage{tabu}
%\usepackage[defaultlines=3,all]{nowidow}
%\usepackage{setspace}
%\usepackage{geometry}
%\usepackage[all]{xy}
%\usepackage{natbib}
%\usepackage{subfigure}
% \usepackage{svg}   % needs --shell-escape + inkscape and is rejected by arXiv; pre-convert SVGs to PDF instead

% ==== math notation ====
% Symbol macros are wrapped in \ensuremath, so they also work in text (\R, \bx, ... auto-enter math).

% ---- quote marks in math ----
\DeclareMathSymbol{\mlq}{\mathord}{operators}{``}
\DeclareMathSymbol{\mrq}{\mathord}{operators}{`'}

% ---- number sets (blackboard bold) ----
% NOTE: \P overrides the built-in pilcrow (¶); see also \S \O \L \H below.
\newcommand{\R}{\ensuremath{\mathbb{R}}}
\newcommand{\N}{\ensuremath{\mathbb{N}}}
\newcommand{\Z}{\ensuremath{\mathbb{Z}}}
\newcommand{\Q}{\ensuremath{\mathbb{Q}}}
\newcommand{\C}{\ensuremath{\mathbb{C}}}
\newcommand{\D}{\ensuremath{\mathbb{D}}}
\newcommand{\E}{\ensuremath{\mathbb{E}}}
\newcommand{\X}{\ensuremath{\mathbb{X}}}
\newcommand{\Y}{\ensuremath{\mathbb{Y}}}
\renewcommand{\P}{\ensuremath{\mathbb{P}}}
\newcommand{\sigB}{\ensuremath{\mathbb{B}}}
\newcommand{\1}{\ensuremath{\mathds{1}}}

% ---- calligraphic sets ----
% NOTE: \S \O \L \H override built-in text commands (§, Ø, Ł, and the o-double-acute accent \H{o}).
\newcommand{\calN}{\ensuremath{\mathcal{N}}}
\newcommand{\calD}{\ensuremath{\mathcal{D}}}
\newcommand{\A}{\ensuremath{\mathcal{A}}}
\newcommand{\B}{\ensuremath{\mathcal{B}}}
\newcommand{\I}{\ensuremath{\mathcal{I}}}
\newcommand{\F}{\ensuremath{\mathcal{F}}}
\newcommand{\G}{\ensuremath{\mathcal{G}}}
\newcommand{\M}{\ensuremath{\mathcal{M}}}
\newcommand{\Tau}{\ensuremath{\mathcal{T}}}
\newcommand{\W}{\ensuremath{\mathcal{W}}}
\newcommand{\sigE}{\ensuremath{\mathcal{E}}}
\renewcommand{\H}{\ensuremath{\mathcal{H}}}
\renewcommand{\L}{\ensuremath{\mathcal{L}}}
\renewcommand{\S}{\ensuremath{\mathcal{S}}}
\renewcommand{\O}{\ensuremath{\mathcal{O}}}

% ---- named operators ----
% NOTE: \Re \Im are redefined to roman "Re"/"Im" (not the default blackletter ℜ/ℑ).
\DeclareMathOperator{\softmax}{softmax}
\DeclareMathOperator*{\argmax}{arg\,max}   % * = subscript sits underneath in display
\DeclareMathOperator{\Var}{Var}
\newcommand{\cov}{\ensuremath{\operatorname{cov}}}
\newcommand{\corr}{\ensuremath{\operatorname{corr}}}
\newcommand{\spanop}{\ensuremath{\operatorname{span}}}
\newcommand{\att}{\ensuremath{\operatorname{attention}}}
\newcommand{\suff}{\ensuremath{\operatorname{suff}}}
\newcommand{\comp}{\ensuremath{\operatorname{comp}}}
\newcommand{\In}{\ensuremath{\operatorname{In}}}
\newcommand{\opcap}{\ensuremath{\operatorname{cap}}}
\newcommand{\entr}{\ensuremath{\operatorname{entries}}}
\newcommand{\ops}{\ensuremath{\operatorname{ops}}}
\newcommand{\id}{\ensuremath{\mathrm{id}}}
\renewcommand{\Re}{\ensuremath{\operatorname{Re}}}
\renewcommand{\Im}{\ensuremath{\operatorname{Im}}}

% ---- symbols & shorthands ----
\newcommand{\eps}{\ensuremath{\varepsilon}}
\renewcommand{\phi}{\ensuremath{\varphi}}     % deliberate: \phi -> varphi
\newcommand{\del}{\ensuremath{\partial}}
\newcommand{\grad}{\ensuremath{\nabla}}
\newcommand{\indep}{\ensuremath{\perp\!\!\!\perp}}
\newcommand{\emptyword}{\ensuremath{\varnothing}}
\newcommand{\tensor}{\ensuremath{\otimes}}
\newcommand{\bigtensor}{\ensuremath{\bigotimes}}
\newcommand{\To}{\ensuremath{\longrightarrow}}
\newcommand{\eqover}[1]{\ensuremath{\stackrel{#1}{=}}}
\newcommand{\darover}[1]{\ensuremath{\xrightarrow{#1}}}
\newcommand{\Hoel}{\text{Höl}}

% ---- paired delimiters ----
\DeclarePairedDelimiter{\norm}{\lVert}{\rVert}   % \norm{x}, \norm*{x} (full size), \norm[\big]{x}
\DeclarePairedDelimiter{\abs}{\lvert}{\rvert}

% ---- bold vectors / matrices ----
\newcommand{\bx}{\ensuremath{\mathbf{x}}}
\newcommand{\bX}{\ensuremath{\mathbf{X}}}
\newcommand{\by}{\ensuremath{\mathbf{y}}}
\newcommand{\bz}{\ensuremath{\mathbf{z}}}
\newcommand{\bB}{\ensuremath{\mathbf{B}}}
\newcommand{\bp}{\ensuremath{\mathbf{p}}}
\newcommand{\bW}{\ensuremath{\mathbf{W}}}
\newcommand{\bt}{\ensuremath{\mathbf{t}}}
\newcommand{\be}{\ensuremath{\mathbf{e}}}

% ---- small constructs ----
\newcommand{\Ball}[2]{\ensuremath{B_#1(#2)}}   % \Ball{\eps}{center}
\newcommand{\twobytwo}[4]{\ensuremath{\begin{pmatrix} #1 & #2 \\ #3 & #4 \end{pmatrix}}}

% ==== text, color & layout helpers ====
\definecolor{darkgreen}{HTML}{009B55}
\newcommand{\gtxt}[1]{\text{\textcolor{gray}{#1}}}
\newcommand{\grntxt}[1]{\text{\textcolor{darkgreen}{#1}}}
\newcommand{\rdtxt}[1]{\text{\textcolor{red}{#1}}}
\newcommand{\code}[1]{\texttt{#1}}
\newcommand{\tldr}{\textbf{TL;DR:}\xspace}
\newcommand{\parahead}[1]{\noindent\textbf{#1.}\enspace}  % run-in heading; use instead of \paragraph to save space
\newcommand*\rot{\rotatebox{90}}

% ---- table cell helpers ----
\newcommand{\cmark}{\ding{51}}
\newcommand{\xmark}{\ding{55}}
\newcommand{\rowhighlight}{\rowcolor{gray!30}}
\newcommand{\cellhighlight}{\cellcolor{yellow!40}}

% ---- rebuttal reviewer tags: bold, distinct colors to signpost which reviewer a point answers ----
\newcommand{\Rone}{\textbf{\textcolor[HTML]{1F77B4}{R1}}}    % blue
\newcommand{\Rtwo}{\textbf{\textcolor[HTML]{D62728}{R2}}}    % red
\newcommand{\Rthree}{\textbf{\textcolor[HTML]{2CA02C}{R3}}}  % green
\newcommand{\Rfour}{\textbf{\textcolor[HTML]{D9730D}{R4}}}   % muted orange

% ---- takeaway callout box ----
\newtcolorbox{takeaway}[1][Takeaway]{
  colback=blue!4!white,      % very light fill
  colframe=blue!50!black,    % muted dark-blue frame
  boxrule=0.6pt,
  arc=2pt,
  left=6pt, right=6pt, top=2pt, bottom=2pt,
  title=#1,
  fonttitle=\bfseries\small,
  coltitle=white,
  colbacktitle=blue!50!black
}

% ==== links & smart references: load LAST (delete if your class already loads hyperref) ====
\usepackage[hidelinks]{hyperref}
\usepackage[capitalize,noabbrev]{cleveref}   % \cref{...}, \Cref{...} — must come after hyperref

% ==== draft mode & \todo ====
% "Explicit" = the main .tex defined \ifdraftmode before % latex-rules — reusable preamble (see the latex-rules skill).
% Reusable macros only; define paper-specific notation in your main .tex, not here.

% ==== packages ====

% ---- fonts & encoding ----
\usepackage[T1]{fontenc}   % proper 8-bit fonts + real hyphenation
\usepackage{lmodern}       % scalable Latin Modern (fixes bitmap CM warnings)

% ---- math ----
\usepackage{amsmath}
\usepackage{amssymb}
\usepackage{mathtools}     % loads + extends amsmath
\usepackage{dsfont}        % \mathds{1}

% ---- graphics & drawing ----
\usepackage{graphicx}
\usepackage[export]{adjustbox}
\usepackage{tikz}
\usetikzlibrary{shapes}
\usetikzlibrary{positioning}
\usetikzlibrary{trees}
%\usetikzlibrary{graphs,graphdrawing}
%\usegdlibrary{trees}
\usepackage[most]{tcolorbox}
\usepackage{subcaption}    % NOT subfigure / subfig
\usepackage{pdflscape}

% ---- color ----
\usepackage{xcolor}

% ---- tables ----
\usepackage{booktabs}
\usepackage{multirow}
\usepackage{makecell}
\usepackage{colortbl}

% ---- typography, lists, units ----
\usepackage[protrusion=true,expansion=true]{microtype}  % char protrusion + font expansion; drop expansion under XeLaTeX (unsupported)
\usepackage{csquotes}
\usepackage{enumitem}
\usepackage{siunitx}
\sisetup{detect-all, per-mode=symbol}   % detect-all: match surrounding font; per-mode=symbol: m/s not m s^-1
\usepackage{xspace}

% ---- floats, algorithms, symbols ----
\usepackage{float}
\usepackage{afterpage}
\usepackage{footmisc}
\usepackage{listings}
\usepackage{algorithm}
\usepackage{algpseudocode}   % algorithmicx; modern replacement for algorithmic
\usepackage{pifont}
%\usepackage{pgfplots}

% ---- occasionally useful (uncomment as needed) ----
%\usepackage{tabu}
%\usepackage[defaultlines=3,all]{nowidow}
%\usepackage{setspace}
%\usepackage{geometry}
%\usepackage[all]{xy}
%\usepackage{natbib}
%\usepackage{subfigure}
% \usepackage{svg}   % needs --shell-escape + inkscape and is rejected by arXiv; pre-convert SVGs to PDF instead

% ==== math notation ====
% Symbol macros are wrapped in \ensuremath, so they also work in text (\R, \bx, ... auto-enter math).

% ---- quote marks in math ----
\DeclareMathSymbol{\mlq}{\mathord}{operators}{``}
\DeclareMathSymbol{\mrq}{\mathord}{operators}{`'}

% ---- number sets (blackboard bold) ----
% NOTE: \P overrides the built-in pilcrow (¶); see also \S \O \L \H below.
\newcommand{\R}{\ensuremath{\mathbb{R}}}
\newcommand{\N}{\ensuremath{\mathbb{N}}}
\newcommand{\Z}{\ensuremath{\mathbb{Z}}}
\newcommand{\Q}{\ensuremath{\mathbb{Q}}}
\newcommand{\C}{\ensuremath{\mathbb{C}}}
\newcommand{\D}{\ensuremath{\mathbb{D}}}
\newcommand{\E}{\ensuremath{\mathbb{E}}}
\newcommand{\X}{\ensuremath{\mathbb{X}}}
\newcommand{\Y}{\ensuremath{\mathbb{Y}}}
\renewcommand{\P}{\ensuremath{\mathbb{P}}}
\newcommand{\sigB}{\ensuremath{\mathbb{B}}}
\newcommand{\1}{\ensuremath{\mathds{1}}}

% ---- calligraphic sets ----
% NOTE: \S \O \L \H override built-in text commands (§, Ø, Ł, and the o-double-acute accent \H{o}).
\newcommand{\calN}{\ensuremath{\mathcal{N}}}
\newcommand{\calD}{\ensuremath{\mathcal{D}}}
\newcommand{\A}{\ensuremath{\mathcal{A}}}
\newcommand{\B}{\ensuremath{\mathcal{B}}}
\newcommand{\I}{\ensuremath{\mathcal{I}}}
\newcommand{\F}{\ensuremath{\mathcal{F}}}
\newcommand{\G}{\ensuremath{\mathcal{G}}}
\newcommand{\M}{\ensuremath{\mathcal{M}}}
\newcommand{\Tau}{\ensuremath{\mathcal{T}}}
\newcommand{\W}{\ensuremath{\mathcal{W}}}
\newcommand{\sigE}{\ensuremath{\mathcal{E}}}
\renewcommand{\H}{\ensuremath{\mathcal{H}}}
\renewcommand{\L}{\ensuremath{\mathcal{L}}}
\renewcommand{\S}{\ensuremath{\mathcal{S}}}
\renewcommand{\O}{\ensuremath{\mathcal{O}}}

% ---- named operators ----
% NOTE: \Re \Im are redefined to roman "Re"/"Im" (not the default blackletter ℜ/ℑ).
\DeclareMathOperator{\softmax}{softmax}
\DeclareMathOperator*{\argmax}{arg\,max}   % * = subscript sits underneath in display
\DeclareMathOperator{\Var}{Var}
\newcommand{\cov}{\ensuremath{\operatorname{cov}}}
\newcommand{\corr}{\ensuremath{\operatorname{corr}}}
\newcommand{\spanop}{\ensuremath{\operatorname{span}}}
\newcommand{\att}{\ensuremath{\operatorname{attention}}}
\newcommand{\suff}{\ensuremath{\operatorname{suff}}}
\newcommand{\comp}{\ensuremath{\operatorname{comp}}}
\newcommand{\In}{\ensuremath{\operatorname{In}}}
\newcommand{\opcap}{\ensuremath{\operatorname{cap}}}
\newcommand{\entr}{\ensuremath{\operatorname{entries}}}
\newcommand{\ops}{\ensuremath{\operatorname{ops}}}
\newcommand{\id}{\ensuremath{\mathrm{id}}}
\renewcommand{\Re}{\ensuremath{\operatorname{Re}}}
\renewcommand{\Im}{\ensuremath{\operatorname{Im}}}

% ---- symbols & shorthands ----
\newcommand{\eps}{\ensuremath{\varepsilon}}
\renewcommand{\phi}{\ensuremath{\varphi}}     % deliberate: \phi -> varphi
\newcommand{\del}{\ensuremath{\partial}}
\newcommand{\grad}{\ensuremath{\nabla}}
\newcommand{\indep}{\ensuremath{\perp\!\!\!\perp}}
\newcommand{\emptyword}{\ensuremath{\varnothing}}
\newcommand{\tensor}{\ensuremath{\otimes}}
\newcommand{\bigtensor}{\ensuremath{\bigotimes}}
\newcommand{\To}{\ensuremath{\longrightarrow}}
\newcommand{\eqover}[1]{\ensuremath{\stackrel{#1}{=}}}
\newcommand{\darover}[1]{\ensuremath{\xrightarrow{#1}}}
\newcommand{\Hoel}{\text{Höl}}

% ---- paired delimiters ----
\DeclarePairedDelimiter{\norm}{\lVert}{\rVert}   % \norm{x}, \norm*{x} (full size), \norm[\big]{x}
\DeclarePairedDelimiter{\abs}{\lvert}{\rvert}

% ---- bold vectors / matrices ----
\newcommand{\bx}{\ensuremath{\mathbf{x}}}
\newcommand{\bX}{\ensuremath{\mathbf{X}}}
\newcommand{\by}{\ensuremath{\mathbf{y}}}
\newcommand{\bz}{\ensuremath{\mathbf{z}}}
\newcommand{\bB}{\ensuremath{\mathbf{B}}}
\newcommand{\bp}{\ensuremath{\mathbf{p}}}
\newcommand{\bW}{\ensuremath{\mathbf{W}}}
\newcommand{\bt}{\ensuremath{\mathbf{t}}}
\newcommand{\be}{\ensuremath{\mathbf{e}}}

% ---- small constructs ----
\newcommand{\Ball}[2]{\ensuremath{B_#1(#2)}}   % \Ball{\eps}{center}
\newcommand{\twobytwo}[4]{\ensuremath{\begin{pmatrix} #1 & #2 \\ #3 & #4 \end{pmatrix}}}

% ==== text, color & layout helpers ====
\definecolor{darkgreen}{HTML}{009B55}
\newcommand{\gtxt}[1]{\text{\textcolor{gray}{#1}}}
\newcommand{\grntxt}[1]{\text{\textcolor{darkgreen}{#1}}}
\newcommand{\rdtxt}[1]{\text{\textcolor{red}{#1}}}
\newcommand{\code}[1]{\texttt{#1}}
\newcommand{\tldr}{\textbf{TL;DR:}\xspace}
\newcommand{\parahead}[1]{\noindent\textbf{#1.}\enspace}  % run-in heading; use instead of \paragraph to save space
\newcommand*\rot{\rotatebox{90}}

% ---- table cell helpers ----
\newcommand{\cmark}{\ding{51}}
\newcommand{\xmark}{\ding{55}}
\newcommand{\rowhighlight}{\rowcolor{gray!30}}
\newcommand{\cellhighlight}{\cellcolor{yellow!40}}

% ---- rebuttal reviewer tags: bold, distinct colors to signpost which reviewer a point answers ----
\newcommand{\Rone}{\textbf{\textcolor[HTML]{1F77B4}{R1}}}    % blue
\newcommand{\Rtwo}{\textbf{\textcolor[HTML]{D62728}{R2}}}    % red
\newcommand{\Rthree}{\textbf{\textcolor[HTML]{2CA02C}{R3}}}  % green
\newcommand{\Rfour}{\textbf{\textcolor[HTML]{D9730D}{R4}}}   % muted orange

% ---- takeaway callout box ----
\newtcolorbox{takeaway}[1][Takeaway]{
  colback=blue!4!white,      % very light fill
  colframe=blue!50!black,    % muted dark-blue frame
  boxrule=0.6pt,
  arc=2pt,
  left=6pt, right=6pt, top=2pt, bottom=2pt,
  title=#1,
  fonttitle=\bfseries\small,
  coltitle=white,
  colbacktitle=blue!50!black
}

% ==== links & smart references: load LAST (delete if your class already loads hyperref) ====
\usepackage[hidelinks]{hyperref}
\usepackage[capitalize,noabbrev]{cleveref}   % \cref{...}, \Cref{...} — must come after hyperref

% ==== draft mode & \todo ====
% "Explicit" = the main .tex defined \ifdraftmode before \input{packages}.
% (The \ifcsname/\csname dance avoids skipping over an already-defined \ifdraftmode, which
%  crashes the naive \ifdefined...\newif idiom.)
\newif\ifdraftmodeexplicit
\ifcsname ifdraftmode\endcsname \draftmodeexplicittrue \else \draftmodeexplicitfalse \fi
\ifdraftmodeexplicit\else \expandafter\newif\csname ifdraftmode\endcsname \fi

% \todo: visible red box by default and in draft mode; a no-op only when draft mode is
% explicitly OFF (\draftmodefalse in the main .tex), so leftover todos vanish in the final build.
\newcommand{\todo}[1]{\colorbox{red}{TODO: #1}}
\ifdraftmodeexplicit \ifdraftmode\else \renewcommand{\todo}[1]{}\fi \fi

% Draft-mode aids (line numbers, label keys, overfull bars) — on only when draft mode is enabled.
% Worth it for long documents / theses; usually overkill for a short single paper.
% Enable by putting "\newif\ifdraftmode \draftmodetrue" in the main .tex BEFORE \input{packages}.
\ifdraftmode
  \usepackage{lineno}\linenumbers   % line numbers for "see line 214" review comments (prose only; align/equation bodies aren't numbered without extra setup)
  \usepackage{showkeys}             % print \label / \ref keys in the margin
  \overfullrule=5pt                 % black bar in the margin marks every overfull line
\fi
.
% (The \ifcsname/\csname dance avoids skipping over an already-defined \ifdraftmode, which
%  crashes the naive \ifdefined...\newif idiom.)
\newif\ifdraftmodeexplicit
\ifcsname ifdraftmode\endcsname \draftmodeexplicittrue \else \draftmodeexplicitfalse \fi
\ifdraftmodeexplicit\else \expandafter\newif\csname ifdraftmode\endcsname \fi

% \todo: visible red box by default and in draft mode; a no-op only when draft mode is
% explicitly OFF (\draftmodefalse in the main .tex), so leftover todos vanish in the final build.
\newcommand{\todo}[1]{\colorbox{red}{TODO: #1}}
\ifdraftmodeexplicit \ifdraftmode\else \renewcommand{\todo}[1]{}\fi \fi

% Draft-mode aids (line numbers, label keys, overfull bars) — on only when draft mode is enabled.
% Worth it for long documents / theses; usually overkill for a short single paper.
% Enable by putting "\newif\ifdraftmode \draftmodetrue" in the main .tex BEFORE % latex-rules — reusable preamble (see the latex-rules skill).
% Reusable macros only; define paper-specific notation in your main .tex, not here.

% ==== packages ====

% ---- fonts & encoding ----
\usepackage[T1]{fontenc}   % proper 8-bit fonts + real hyphenation
\usepackage{lmodern}       % scalable Latin Modern (fixes bitmap CM warnings)

% ---- math ----
\usepackage{amsmath}
\usepackage{amssymb}
\usepackage{mathtools}     % loads + extends amsmath
\usepackage{dsfont}        % \mathds{1}

% ---- graphics & drawing ----
\usepackage{graphicx}
\usepackage[export]{adjustbox}
\usepackage{tikz}
\usetikzlibrary{shapes}
\usetikzlibrary{positioning}
\usetikzlibrary{trees}
%\usetikzlibrary{graphs,graphdrawing}
%\usegdlibrary{trees}
\usepackage[most]{tcolorbox}
\usepackage{subcaption}    % NOT subfigure / subfig
\usepackage{pdflscape}

% ---- color ----
\usepackage{xcolor}

% ---- tables ----
\usepackage{booktabs}
\usepackage{multirow}
\usepackage{makecell}
\usepackage{colortbl}

% ---- typography, lists, units ----
\usepackage[protrusion=true,expansion=true]{microtype}  % char protrusion + font expansion; drop expansion under XeLaTeX (unsupported)
\usepackage{csquotes}
\usepackage{enumitem}
\usepackage{siunitx}
\sisetup{detect-all, per-mode=symbol}   % detect-all: match surrounding font; per-mode=symbol: m/s not m s^-1
\usepackage{xspace}

% ---- floats, algorithms, symbols ----
\usepackage{float}
\usepackage{afterpage}
\usepackage{footmisc}
\usepackage{listings}
\usepackage{algorithm}
\usepackage{algpseudocode}   % algorithmicx; modern replacement for algorithmic
\usepackage{pifont}
%\usepackage{pgfplots}

% ---- occasionally useful (uncomment as needed) ----
%\usepackage{tabu}
%\usepackage[defaultlines=3,all]{nowidow}
%\usepackage{setspace}
%\usepackage{geometry}
%\usepackage[all]{xy}
%\usepackage{natbib}
%\usepackage{subfigure}
% \usepackage{svg}   % needs --shell-escape + inkscape and is rejected by arXiv; pre-convert SVGs to PDF instead

% ==== math notation ====
% Symbol macros are wrapped in \ensuremath, so they also work in text (\R, \bx, ... auto-enter math).

% ---- quote marks in math ----
\DeclareMathSymbol{\mlq}{\mathord}{operators}{``}
\DeclareMathSymbol{\mrq}{\mathord}{operators}{`'}

% ---- number sets (blackboard bold) ----
% NOTE: \P overrides the built-in pilcrow (¶); see also \S \O \L \H below.
\newcommand{\R}{\ensuremath{\mathbb{R}}}
\newcommand{\N}{\ensuremath{\mathbb{N}}}
\newcommand{\Z}{\ensuremath{\mathbb{Z}}}
\newcommand{\Q}{\ensuremath{\mathbb{Q}}}
\newcommand{\C}{\ensuremath{\mathbb{C}}}
\newcommand{\D}{\ensuremath{\mathbb{D}}}
\newcommand{\E}{\ensuremath{\mathbb{E}}}
\newcommand{\X}{\ensuremath{\mathbb{X}}}
\newcommand{\Y}{\ensuremath{\mathbb{Y}}}
\renewcommand{\P}{\ensuremath{\mathbb{P}}}
\newcommand{\sigB}{\ensuremath{\mathbb{B}}}
\newcommand{\1}{\ensuremath{\mathds{1}}}

% ---- calligraphic sets ----
% NOTE: \S \O \L \H override built-in text commands (§, Ø, Ł, and the o-double-acute accent \H{o}).
\newcommand{\calN}{\ensuremath{\mathcal{N}}}
\newcommand{\calD}{\ensuremath{\mathcal{D}}}
\newcommand{\A}{\ensuremath{\mathcal{A}}}
\newcommand{\B}{\ensuremath{\mathcal{B}}}
\newcommand{\I}{\ensuremath{\mathcal{I}}}
\newcommand{\F}{\ensuremath{\mathcal{F}}}
\newcommand{\G}{\ensuremath{\mathcal{G}}}
\newcommand{\M}{\ensuremath{\mathcal{M}}}
\newcommand{\Tau}{\ensuremath{\mathcal{T}}}
\newcommand{\W}{\ensuremath{\mathcal{W}}}
\newcommand{\sigE}{\ensuremath{\mathcal{E}}}
\renewcommand{\H}{\ensuremath{\mathcal{H}}}
\renewcommand{\L}{\ensuremath{\mathcal{L}}}
\renewcommand{\S}{\ensuremath{\mathcal{S}}}
\renewcommand{\O}{\ensuremath{\mathcal{O}}}

% ---- named operators ----
% NOTE: \Re \Im are redefined to roman "Re"/"Im" (not the default blackletter ℜ/ℑ).
\DeclareMathOperator{\softmax}{softmax}
\DeclareMathOperator*{\argmax}{arg\,max}   % * = subscript sits underneath in display
\DeclareMathOperator{\Var}{Var}
\newcommand{\cov}{\ensuremath{\operatorname{cov}}}
\newcommand{\corr}{\ensuremath{\operatorname{corr}}}
\newcommand{\spanop}{\ensuremath{\operatorname{span}}}
\newcommand{\att}{\ensuremath{\operatorname{attention}}}
\newcommand{\suff}{\ensuremath{\operatorname{suff}}}
\newcommand{\comp}{\ensuremath{\operatorname{comp}}}
\newcommand{\In}{\ensuremath{\operatorname{In}}}
\newcommand{\opcap}{\ensuremath{\operatorname{cap}}}
\newcommand{\entr}{\ensuremath{\operatorname{entries}}}
\newcommand{\ops}{\ensuremath{\operatorname{ops}}}
\newcommand{\id}{\ensuremath{\mathrm{id}}}
\renewcommand{\Re}{\ensuremath{\operatorname{Re}}}
\renewcommand{\Im}{\ensuremath{\operatorname{Im}}}

% ---- symbols & shorthands ----
\newcommand{\eps}{\ensuremath{\varepsilon}}
\renewcommand{\phi}{\ensuremath{\varphi}}     % deliberate: \phi -> varphi
\newcommand{\del}{\ensuremath{\partial}}
\newcommand{\grad}{\ensuremath{\nabla}}
\newcommand{\indep}{\ensuremath{\perp\!\!\!\perp}}
\newcommand{\emptyword}{\ensuremath{\varnothing}}
\newcommand{\tensor}{\ensuremath{\otimes}}
\newcommand{\bigtensor}{\ensuremath{\bigotimes}}
\newcommand{\To}{\ensuremath{\longrightarrow}}
\newcommand{\eqover}[1]{\ensuremath{\stackrel{#1}{=}}}
\newcommand{\darover}[1]{\ensuremath{\xrightarrow{#1}}}
\newcommand{\Hoel}{\text{Höl}}

% ---- paired delimiters ----
\DeclarePairedDelimiter{\norm}{\lVert}{\rVert}   % \norm{x}, \norm*{x} (full size), \norm[\big]{x}
\DeclarePairedDelimiter{\abs}{\lvert}{\rvert}

% ---- bold vectors / matrices ----
\newcommand{\bx}{\ensuremath{\mathbf{x}}}
\newcommand{\bX}{\ensuremath{\mathbf{X}}}
\newcommand{\by}{\ensuremath{\mathbf{y}}}
\newcommand{\bz}{\ensuremath{\mathbf{z}}}
\newcommand{\bB}{\ensuremath{\mathbf{B}}}
\newcommand{\bp}{\ensuremath{\mathbf{p}}}
\newcommand{\bW}{\ensuremath{\mathbf{W}}}
\newcommand{\bt}{\ensuremath{\mathbf{t}}}
\newcommand{\be}{\ensuremath{\mathbf{e}}}

% ---- small constructs ----
\newcommand{\Ball}[2]{\ensuremath{B_#1(#2)}}   % \Ball{\eps}{center}
\newcommand{\twobytwo}[4]{\ensuremath{\begin{pmatrix} #1 & #2 \\ #3 & #4 \end{pmatrix}}}

% ==== text, color & layout helpers ====
\definecolor{darkgreen}{HTML}{009B55}
\newcommand{\gtxt}[1]{\text{\textcolor{gray}{#1}}}
\newcommand{\grntxt}[1]{\text{\textcolor{darkgreen}{#1}}}
\newcommand{\rdtxt}[1]{\text{\textcolor{red}{#1}}}
\newcommand{\code}[1]{\texttt{#1}}
\newcommand{\tldr}{\textbf{TL;DR:}\xspace}
\newcommand{\parahead}[1]{\noindent\textbf{#1.}\enspace}  % run-in heading; use instead of \paragraph to save space
\newcommand*\rot{\rotatebox{90}}

% ---- table cell helpers ----
\newcommand{\cmark}{\ding{51}}
\newcommand{\xmark}{\ding{55}}
\newcommand{\rowhighlight}{\rowcolor{gray!30}}
\newcommand{\cellhighlight}{\cellcolor{yellow!40}}

% ---- rebuttal reviewer tags: bold, distinct colors to signpost which reviewer a point answers ----
\newcommand{\Rone}{\textbf{\textcolor[HTML]{1F77B4}{R1}}}    % blue
\newcommand{\Rtwo}{\textbf{\textcolor[HTML]{D62728}{R2}}}    % red
\newcommand{\Rthree}{\textbf{\textcolor[HTML]{2CA02C}{R3}}}  % green
\newcommand{\Rfour}{\textbf{\textcolor[HTML]{D9730D}{R4}}}   % muted orange

% ---- takeaway callout box ----
\newtcolorbox{takeaway}[1][Takeaway]{
  colback=blue!4!white,      % very light fill
  colframe=blue!50!black,    % muted dark-blue frame
  boxrule=0.6pt,
  arc=2pt,
  left=6pt, right=6pt, top=2pt, bottom=2pt,
  title=#1,
  fonttitle=\bfseries\small,
  coltitle=white,
  colbacktitle=blue!50!black
}

% ==== links & smart references: load LAST (delete if your class already loads hyperref) ====
\usepackage[hidelinks]{hyperref}
\usepackage[capitalize,noabbrev]{cleveref}   % \cref{...}, \Cref{...} — must come after hyperref

% ==== draft mode & \todo ====
% "Explicit" = the main .tex defined \ifdraftmode before \input{packages}.
% (The \ifcsname/\csname dance avoids skipping over an already-defined \ifdraftmode, which
%  crashes the naive \ifdefined...\newif idiom.)
\newif\ifdraftmodeexplicit
\ifcsname ifdraftmode\endcsname \draftmodeexplicittrue \else \draftmodeexplicitfalse \fi
\ifdraftmodeexplicit\else \expandafter\newif\csname ifdraftmode\endcsname \fi

% \todo: visible red box by default and in draft mode; a no-op only when draft mode is
% explicitly OFF (\draftmodefalse in the main .tex), so leftover todos vanish in the final build.
\newcommand{\todo}[1]{\colorbox{red}{TODO: #1}}
\ifdraftmodeexplicit \ifdraftmode\else \renewcommand{\todo}[1]{}\fi \fi

% Draft-mode aids (line numbers, label keys, overfull bars) — on only when draft mode is enabled.
% Worth it for long documents / theses; usually overkill for a short single paper.
% Enable by putting "\newif\ifdraftmode \draftmodetrue" in the main .tex BEFORE \input{packages}.
\ifdraftmode
  \usepackage{lineno}\linenumbers   % line numbers for "see line 214" review comments (prose only; align/equation bodies aren't numbered without extra setup)
  \usepackage{showkeys}             % print \label / \ref keys in the margin
  \overfullrule=5pt                 % black bar in the margin marks every overfull line
\fi
.
\ifdraftmode
  \usepackage{lineno}\linenumbers   % line numbers for "see line 214" review comments (prose only; align/equation bodies aren't numbered without extra setup)
  \usepackage{showkeys}             % print \label / \ref keys in the margin
  \overfullrule=5pt                 % black bar in the margin marks every overfull line
\fi
.
% (The \ifcsname/\csname dance avoids skipping over an already-defined \ifdraftmode, which
%  crashes the naive \ifdefined...\newif idiom.)
\newif\ifdraftmodeexplicit
\ifcsname ifdraftmode\endcsname \draftmodeexplicittrue \else \draftmodeexplicitfalse \fi
\ifdraftmodeexplicit\else \expandafter\newif\csname ifdraftmode\endcsname \fi

% \todo: visible red box by default and in draft mode; a no-op only when draft mode is
% explicitly OFF (\draftmodefalse in the main .tex), so leftover todos vanish in the final build.
\newcommand{\todo}[1]{\colorbox{red}{TODO: #1}}
\ifdraftmodeexplicit \ifdraftmode\else \renewcommand{\todo}[1]{}\fi \fi

% Draft-mode aids (line numbers, label keys, overfull bars) — on only when draft mode is enabled.
% Worth it for long documents / theses; usually overkill for a short single paper.
% Enable by putting "\newif\ifdraftmode \draftmodetrue" in the main .tex BEFORE % latex-rules — reusable preamble (see the latex-rules skill).
% Reusable macros only; define paper-specific notation in your main .tex, not here.

% ==== packages ====

% ---- fonts & encoding ----
\usepackage[T1]{fontenc}   % proper 8-bit fonts + real hyphenation
\usepackage{lmodern}       % scalable Latin Modern (fixes bitmap CM warnings)

% ---- math ----
\usepackage{amsmath}
\usepackage{amssymb}
\usepackage{mathtools}     % loads + extends amsmath
\usepackage{dsfont}        % \mathds{1}

% ---- graphics & drawing ----
\usepackage{graphicx}
\usepackage[export]{adjustbox}
\usepackage{tikz}
\usetikzlibrary{shapes}
\usetikzlibrary{positioning}
\usetikzlibrary{trees}
%\usetikzlibrary{graphs,graphdrawing}
%\usegdlibrary{trees}
\usepackage[most]{tcolorbox}
\usepackage{subcaption}    % NOT subfigure / subfig
\usepackage{pdflscape}

% ---- color ----
\usepackage{xcolor}

% ---- tables ----
\usepackage{booktabs}
\usepackage{multirow}
\usepackage{makecell}
\usepackage{colortbl}

% ---- typography, lists, units ----
\usepackage[protrusion=true,expansion=true]{microtype}  % char protrusion + font expansion; drop expansion under XeLaTeX (unsupported)
\usepackage{csquotes}
\usepackage{enumitem}
\usepackage{siunitx}
\sisetup{detect-all, per-mode=symbol}   % detect-all: match surrounding font; per-mode=symbol: m/s not m s^-1
\usepackage{xspace}

% ---- floats, algorithms, symbols ----
\usepackage{float}
\usepackage{afterpage}
\usepackage{footmisc}
\usepackage{listings}
\usepackage{algorithm}
\usepackage{algpseudocode}   % algorithmicx; modern replacement for algorithmic
\usepackage{pifont}
%\usepackage{pgfplots}

% ---- occasionally useful (uncomment as needed) ----
%\usepackage{tabu}
%\usepackage[defaultlines=3,all]{nowidow}
%\usepackage{setspace}
%\usepackage{geometry}
%\usepackage[all]{xy}
%\usepackage{natbib}
%\usepackage{subfigure}
% \usepackage{svg}   % needs --shell-escape + inkscape and is rejected by arXiv; pre-convert SVGs to PDF instead

% ==== math notation ====
% Symbol macros are wrapped in \ensuremath, so they also work in text (\R, \bx, ... auto-enter math).

% ---- quote marks in math ----
\DeclareMathSymbol{\mlq}{\mathord}{operators}{``}
\DeclareMathSymbol{\mrq}{\mathord}{operators}{`'}

% ---- number sets (blackboard bold) ----
% NOTE: \P overrides the built-in pilcrow (¶); see also \S \O \L \H below.
\newcommand{\R}{\ensuremath{\mathbb{R}}}
\newcommand{\N}{\ensuremath{\mathbb{N}}}
\newcommand{\Z}{\ensuremath{\mathbb{Z}}}
\newcommand{\Q}{\ensuremath{\mathbb{Q}}}
\newcommand{\C}{\ensuremath{\mathbb{C}}}
\newcommand{\D}{\ensuremath{\mathbb{D}}}
\newcommand{\E}{\ensuremath{\mathbb{E}}}
\newcommand{\X}{\ensuremath{\mathbb{X}}}
\newcommand{\Y}{\ensuremath{\mathbb{Y}}}
\renewcommand{\P}{\ensuremath{\mathbb{P}}}
\newcommand{\sigB}{\ensuremath{\mathbb{B}}}
\newcommand{\1}{\ensuremath{\mathds{1}}}

% ---- calligraphic sets ----
% NOTE: \S \O \L \H override built-in text commands (§, Ø, Ł, and the o-double-acute accent \H{o}).
\newcommand{\calN}{\ensuremath{\mathcal{N}}}
\newcommand{\calD}{\ensuremath{\mathcal{D}}}
\newcommand{\A}{\ensuremath{\mathcal{A}}}
\newcommand{\B}{\ensuremath{\mathcal{B}}}
\newcommand{\I}{\ensuremath{\mathcal{I}}}
\newcommand{\F}{\ensuremath{\mathcal{F}}}
\newcommand{\G}{\ensuremath{\mathcal{G}}}
\newcommand{\M}{\ensuremath{\mathcal{M}}}
\newcommand{\Tau}{\ensuremath{\mathcal{T}}}
\newcommand{\W}{\ensuremath{\mathcal{W}}}
\newcommand{\sigE}{\ensuremath{\mathcal{E}}}
\renewcommand{\H}{\ensuremath{\mathcal{H}}}
\renewcommand{\L}{\ensuremath{\mathcal{L}}}
\renewcommand{\S}{\ensuremath{\mathcal{S}}}
\renewcommand{\O}{\ensuremath{\mathcal{O}}}

% ---- named operators ----
% NOTE: \Re \Im are redefined to roman "Re"/"Im" (not the default blackletter ℜ/ℑ).
\DeclareMathOperator{\softmax}{softmax}
\DeclareMathOperator*{\argmax}{arg\,max}   % * = subscript sits underneath in display
\DeclareMathOperator{\Var}{Var}
\newcommand{\cov}{\ensuremath{\operatorname{cov}}}
\newcommand{\corr}{\ensuremath{\operatorname{corr}}}
\newcommand{\spanop}{\ensuremath{\operatorname{span}}}
\newcommand{\att}{\ensuremath{\operatorname{attention}}}
\newcommand{\suff}{\ensuremath{\operatorname{suff}}}
\newcommand{\comp}{\ensuremath{\operatorname{comp}}}
\newcommand{\In}{\ensuremath{\operatorname{In}}}
\newcommand{\opcap}{\ensuremath{\operatorname{cap}}}
\newcommand{\entr}{\ensuremath{\operatorname{entries}}}
\newcommand{\ops}{\ensuremath{\operatorname{ops}}}
\newcommand{\id}{\ensuremath{\mathrm{id}}}
\renewcommand{\Re}{\ensuremath{\operatorname{Re}}}
\renewcommand{\Im}{\ensuremath{\operatorname{Im}}}

% ---- symbols & shorthands ----
\newcommand{\eps}{\ensuremath{\varepsilon}}
\renewcommand{\phi}{\ensuremath{\varphi}}     % deliberate: \phi -> varphi
\newcommand{\del}{\ensuremath{\partial}}
\newcommand{\grad}{\ensuremath{\nabla}}
\newcommand{\indep}{\ensuremath{\perp\!\!\!\perp}}
\newcommand{\emptyword}{\ensuremath{\varnothing}}
\newcommand{\tensor}{\ensuremath{\otimes}}
\newcommand{\bigtensor}{\ensuremath{\bigotimes}}
\newcommand{\To}{\ensuremath{\longrightarrow}}
\newcommand{\eqover}[1]{\ensuremath{\stackrel{#1}{=}}}
\newcommand{\darover}[1]{\ensuremath{\xrightarrow{#1}}}
\newcommand{\Hoel}{\text{Höl}}

% ---- paired delimiters ----
\DeclarePairedDelimiter{\norm}{\lVert}{\rVert}   % \norm{x}, \norm*{x} (full size), \norm[\big]{x}
\DeclarePairedDelimiter{\abs}{\lvert}{\rvert}

% ---- bold vectors / matrices ----
\newcommand{\bx}{\ensuremath{\mathbf{x}}}
\newcommand{\bX}{\ensuremath{\mathbf{X}}}
\newcommand{\by}{\ensuremath{\mathbf{y}}}
\newcommand{\bz}{\ensuremath{\mathbf{z}}}
\newcommand{\bB}{\ensuremath{\mathbf{B}}}
\newcommand{\bp}{\ensuremath{\mathbf{p}}}
\newcommand{\bW}{\ensuremath{\mathbf{W}}}
\newcommand{\bt}{\ensuremath{\mathbf{t}}}
\newcommand{\be}{\ensuremath{\mathbf{e}}}

% ---- small constructs ----
\newcommand{\Ball}[2]{\ensuremath{B_#1(#2)}}   % \Ball{\eps}{center}
\newcommand{\twobytwo}[4]{\ensuremath{\begin{pmatrix} #1 & #2 \\ #3 & #4 \end{pmatrix}}}

% ==== text, color & layout helpers ====
\definecolor{darkgreen}{HTML}{009B55}
\newcommand{\gtxt}[1]{\text{\textcolor{gray}{#1}}}
\newcommand{\grntxt}[1]{\text{\textcolor{darkgreen}{#1}}}
\newcommand{\rdtxt}[1]{\text{\textcolor{red}{#1}}}
\newcommand{\code}[1]{\texttt{#1}}
\newcommand{\tldr}{\textbf{TL;DR:}\xspace}
\newcommand{\parahead}[1]{\noindent\textbf{#1.}\enspace}  % run-in heading; use instead of \paragraph to save space
\newcommand*\rot{\rotatebox{90}}

% ---- table cell helpers ----
\newcommand{\cmark}{\ding{51}}
\newcommand{\xmark}{\ding{55}}
\newcommand{\rowhighlight}{\rowcolor{gray!30}}
\newcommand{\cellhighlight}{\cellcolor{yellow!40}}

% ---- rebuttal reviewer tags: bold, distinct colors to signpost which reviewer a point answers ----
\newcommand{\Rone}{\textbf{\textcolor[HTML]{1F77B4}{R1}}}    % blue
\newcommand{\Rtwo}{\textbf{\textcolor[HTML]{D62728}{R2}}}    % red
\newcommand{\Rthree}{\textbf{\textcolor[HTML]{2CA02C}{R3}}}  % green
\newcommand{\Rfour}{\textbf{\textcolor[HTML]{D9730D}{R4}}}   % muted orange

% ---- takeaway callout box ----
\newtcolorbox{takeaway}[1][Takeaway]{
  colback=blue!4!white,      % very light fill
  colframe=blue!50!black,    % muted dark-blue frame
  boxrule=0.6pt,
  arc=2pt,
  left=6pt, right=6pt, top=2pt, bottom=2pt,
  title=#1,
  fonttitle=\bfseries\small,
  coltitle=white,
  colbacktitle=blue!50!black
}

% ==== links & smart references: load LAST (delete if your class already loads hyperref) ====
\usepackage[hidelinks]{hyperref}
\usepackage[capitalize,noabbrev]{cleveref}   % \cref{...}, \Cref{...} — must come after hyperref

% ==== draft mode & \todo ====
% "Explicit" = the main .tex defined \ifdraftmode before % latex-rules — reusable preamble (see the latex-rules skill).
% Reusable macros only; define paper-specific notation in your main .tex, not here.

% ==== packages ====

% ---- fonts & encoding ----
\usepackage[T1]{fontenc}   % proper 8-bit fonts + real hyphenation
\usepackage{lmodern}       % scalable Latin Modern (fixes bitmap CM warnings)

% ---- math ----
\usepackage{amsmath}
\usepackage{amssymb}
\usepackage{mathtools}     % loads + extends amsmath
\usepackage{dsfont}        % \mathds{1}

% ---- graphics & drawing ----
\usepackage{graphicx}
\usepackage[export]{adjustbox}
\usepackage{tikz}
\usetikzlibrary{shapes}
\usetikzlibrary{positioning}
\usetikzlibrary{trees}
%\usetikzlibrary{graphs,graphdrawing}
%\usegdlibrary{trees}
\usepackage[most]{tcolorbox}
\usepackage{subcaption}    % NOT subfigure / subfig
\usepackage{pdflscape}

% ---- color ----
\usepackage{xcolor}

% ---- tables ----
\usepackage{booktabs}
\usepackage{multirow}
\usepackage{makecell}
\usepackage{colortbl}

% ---- typography, lists, units ----
\usepackage[protrusion=true,expansion=true]{microtype}  % char protrusion + font expansion; drop expansion under XeLaTeX (unsupported)
\usepackage{csquotes}
\usepackage{enumitem}
\usepackage{siunitx}
\sisetup{detect-all, per-mode=symbol}   % detect-all: match surrounding font; per-mode=symbol: m/s not m s^-1
\usepackage{xspace}

% ---- floats, algorithms, symbols ----
\usepackage{float}
\usepackage{afterpage}
\usepackage{footmisc}
\usepackage{listings}
\usepackage{algorithm}
\usepackage{algpseudocode}   % algorithmicx; modern replacement for algorithmic
\usepackage{pifont}
%\usepackage{pgfplots}

% ---- occasionally useful (uncomment as needed) ----
%\usepackage{tabu}
%\usepackage[defaultlines=3,all]{nowidow}
%\usepackage{setspace}
%\usepackage{geometry}
%\usepackage[all]{xy}
%\usepackage{natbib}
%\usepackage{subfigure}
% \usepackage{svg}   % needs --shell-escape + inkscape and is rejected by arXiv; pre-convert SVGs to PDF instead

% ==== math notation ====
% Symbol macros are wrapped in \ensuremath, so they also work in text (\R, \bx, ... auto-enter math).

% ---- quote marks in math ----
\DeclareMathSymbol{\mlq}{\mathord}{operators}{``}
\DeclareMathSymbol{\mrq}{\mathord}{operators}{`'}

% ---- number sets (blackboard bold) ----
% NOTE: \P overrides the built-in pilcrow (¶); see also \S \O \L \H below.
\newcommand{\R}{\ensuremath{\mathbb{R}}}
\newcommand{\N}{\ensuremath{\mathbb{N}}}
\newcommand{\Z}{\ensuremath{\mathbb{Z}}}
\newcommand{\Q}{\ensuremath{\mathbb{Q}}}
\newcommand{\C}{\ensuremath{\mathbb{C}}}
\newcommand{\D}{\ensuremath{\mathbb{D}}}
\newcommand{\E}{\ensuremath{\mathbb{E}}}
\newcommand{\X}{\ensuremath{\mathbb{X}}}
\newcommand{\Y}{\ensuremath{\mathbb{Y}}}
\renewcommand{\P}{\ensuremath{\mathbb{P}}}
\newcommand{\sigB}{\ensuremath{\mathbb{B}}}
\newcommand{\1}{\ensuremath{\mathds{1}}}

% ---- calligraphic sets ----
% NOTE: \S \O \L \H override built-in text commands (§, Ø, Ł, and the o-double-acute accent \H{o}).
\newcommand{\calN}{\ensuremath{\mathcal{N}}}
\newcommand{\calD}{\ensuremath{\mathcal{D}}}
\newcommand{\A}{\ensuremath{\mathcal{A}}}
\newcommand{\B}{\ensuremath{\mathcal{B}}}
\newcommand{\I}{\ensuremath{\mathcal{I}}}
\newcommand{\F}{\ensuremath{\mathcal{F}}}
\newcommand{\G}{\ensuremath{\mathcal{G}}}
\newcommand{\M}{\ensuremath{\mathcal{M}}}
\newcommand{\Tau}{\ensuremath{\mathcal{T}}}
\newcommand{\W}{\ensuremath{\mathcal{W}}}
\newcommand{\sigE}{\ensuremath{\mathcal{E}}}
\renewcommand{\H}{\ensuremath{\mathcal{H}}}
\renewcommand{\L}{\ensuremath{\mathcal{L}}}
\renewcommand{\S}{\ensuremath{\mathcal{S}}}
\renewcommand{\O}{\ensuremath{\mathcal{O}}}

% ---- named operators ----
% NOTE: \Re \Im are redefined to roman "Re"/"Im" (not the default blackletter ℜ/ℑ).
\DeclareMathOperator{\softmax}{softmax}
\DeclareMathOperator*{\argmax}{arg\,max}   % * = subscript sits underneath in display
\DeclareMathOperator{\Var}{Var}
\newcommand{\cov}{\ensuremath{\operatorname{cov}}}
\newcommand{\corr}{\ensuremath{\operatorname{corr}}}
\newcommand{\spanop}{\ensuremath{\operatorname{span}}}
\newcommand{\att}{\ensuremath{\operatorname{attention}}}
\newcommand{\suff}{\ensuremath{\operatorname{suff}}}
\newcommand{\comp}{\ensuremath{\operatorname{comp}}}
\newcommand{\In}{\ensuremath{\operatorname{In}}}
\newcommand{\opcap}{\ensuremath{\operatorname{cap}}}
\newcommand{\entr}{\ensuremath{\operatorname{entries}}}
\newcommand{\ops}{\ensuremath{\operatorname{ops}}}
\newcommand{\id}{\ensuremath{\mathrm{id}}}
\renewcommand{\Re}{\ensuremath{\operatorname{Re}}}
\renewcommand{\Im}{\ensuremath{\operatorname{Im}}}

% ---- symbols & shorthands ----
\newcommand{\eps}{\ensuremath{\varepsilon}}
\renewcommand{\phi}{\ensuremath{\varphi}}     % deliberate: \phi -> varphi
\newcommand{\del}{\ensuremath{\partial}}
\newcommand{\grad}{\ensuremath{\nabla}}
\newcommand{\indep}{\ensuremath{\perp\!\!\!\perp}}
\newcommand{\emptyword}{\ensuremath{\varnothing}}
\newcommand{\tensor}{\ensuremath{\otimes}}
\newcommand{\bigtensor}{\ensuremath{\bigotimes}}
\newcommand{\To}{\ensuremath{\longrightarrow}}
\newcommand{\eqover}[1]{\ensuremath{\stackrel{#1}{=}}}
\newcommand{\darover}[1]{\ensuremath{\xrightarrow{#1}}}
\newcommand{\Hoel}{\text{Höl}}

% ---- paired delimiters ----
\DeclarePairedDelimiter{\norm}{\lVert}{\rVert}   % \norm{x}, \norm*{x} (full size), \norm[\big]{x}
\DeclarePairedDelimiter{\abs}{\lvert}{\rvert}

% ---- bold vectors / matrices ----
\newcommand{\bx}{\ensuremath{\mathbf{x}}}
\newcommand{\bX}{\ensuremath{\mathbf{X}}}
\newcommand{\by}{\ensuremath{\mathbf{y}}}
\newcommand{\bz}{\ensuremath{\mathbf{z}}}
\newcommand{\bB}{\ensuremath{\mathbf{B}}}
\newcommand{\bp}{\ensuremath{\mathbf{p}}}
\newcommand{\bW}{\ensuremath{\mathbf{W}}}
\newcommand{\bt}{\ensuremath{\mathbf{t}}}
\newcommand{\be}{\ensuremath{\mathbf{e}}}

% ---- small constructs ----
\newcommand{\Ball}[2]{\ensuremath{B_#1(#2)}}   % \Ball{\eps}{center}
\newcommand{\twobytwo}[4]{\ensuremath{\begin{pmatrix} #1 & #2 \\ #3 & #4 \end{pmatrix}}}

% ==== text, color & layout helpers ====
\definecolor{darkgreen}{HTML}{009B55}
\newcommand{\gtxt}[1]{\text{\textcolor{gray}{#1}}}
\newcommand{\grntxt}[1]{\text{\textcolor{darkgreen}{#1}}}
\newcommand{\rdtxt}[1]{\text{\textcolor{red}{#1}}}
\newcommand{\code}[1]{\texttt{#1}}
\newcommand{\tldr}{\textbf{TL;DR:}\xspace}
\newcommand{\parahead}[1]{\noindent\textbf{#1.}\enspace}  % run-in heading; use instead of \paragraph to save space
\newcommand*\rot{\rotatebox{90}}

% ---- table cell helpers ----
\newcommand{\cmark}{\ding{51}}
\newcommand{\xmark}{\ding{55}}
\newcommand{\rowhighlight}{\rowcolor{gray!30}}
\newcommand{\cellhighlight}{\cellcolor{yellow!40}}

% ---- rebuttal reviewer tags: bold, distinct colors to signpost which reviewer a point answers ----
\newcommand{\Rone}{\textbf{\textcolor[HTML]{1F77B4}{R1}}}    % blue
\newcommand{\Rtwo}{\textbf{\textcolor[HTML]{D62728}{R2}}}    % red
\newcommand{\Rthree}{\textbf{\textcolor[HTML]{2CA02C}{R3}}}  % green
\newcommand{\Rfour}{\textbf{\textcolor[HTML]{D9730D}{R4}}}   % muted orange

% ---- takeaway callout box ----
\newtcolorbox{takeaway}[1][Takeaway]{
  colback=blue!4!white,      % very light fill
  colframe=blue!50!black,    % muted dark-blue frame
  boxrule=0.6pt,
  arc=2pt,
  left=6pt, right=6pt, top=2pt, bottom=2pt,
  title=#1,
  fonttitle=\bfseries\small,
  coltitle=white,
  colbacktitle=blue!50!black
}

% ==== links & smart references: load LAST (delete if your class already loads hyperref) ====
\usepackage[hidelinks]{hyperref}
\usepackage[capitalize,noabbrev]{cleveref}   % \cref{...}, \Cref{...} — must come after hyperref

% ==== draft mode & \todo ====
% "Explicit" = the main .tex defined \ifdraftmode before \input{packages}.
% (The \ifcsname/\csname dance avoids skipping over an already-defined \ifdraftmode, which
%  crashes the naive \ifdefined...\newif idiom.)
\newif\ifdraftmodeexplicit
\ifcsname ifdraftmode\endcsname \draftmodeexplicittrue \else \draftmodeexplicitfalse \fi
\ifdraftmodeexplicit\else \expandafter\newif\csname ifdraftmode\endcsname \fi

% \todo: visible red box by default and in draft mode; a no-op only when draft mode is
% explicitly OFF (\draftmodefalse in the main .tex), so leftover todos vanish in the final build.
\newcommand{\todo}[1]{\colorbox{red}{TODO: #1}}
\ifdraftmodeexplicit \ifdraftmode\else \renewcommand{\todo}[1]{}\fi \fi

% Draft-mode aids (line numbers, label keys, overfull bars) — on only when draft mode is enabled.
% Worth it for long documents / theses; usually overkill for a short single paper.
% Enable by putting "\newif\ifdraftmode \draftmodetrue" in the main .tex BEFORE \input{packages}.
\ifdraftmode
  \usepackage{lineno}\linenumbers   % line numbers for "see line 214" review comments (prose only; align/equation bodies aren't numbered without extra setup)
  \usepackage{showkeys}             % print \label / \ref keys in the margin
  \overfullrule=5pt                 % black bar in the margin marks every overfull line
\fi
.
% (The \ifcsname/\csname dance avoids skipping over an already-defined \ifdraftmode, which
%  crashes the naive \ifdefined...\newif idiom.)
\newif\ifdraftmodeexplicit
\ifcsname ifdraftmode\endcsname \draftmodeexplicittrue \else \draftmodeexplicitfalse \fi
\ifdraftmodeexplicit\else \expandafter\newif\csname ifdraftmode\endcsname \fi

% \todo: visible red box by default and in draft mode; a no-op only when draft mode is
% explicitly OFF (\draftmodefalse in the main .tex), so leftover todos vanish in the final build.
\newcommand{\todo}[1]{\colorbox{red}{TODO: #1}}
\ifdraftmodeexplicit \ifdraftmode\else \renewcommand{\todo}[1]{}\fi \fi

% Draft-mode aids (line numbers, label keys, overfull bars) — on only when draft mode is enabled.
% Worth it for long documents / theses; usually overkill for a short single paper.
% Enable by putting "\newif\ifdraftmode \draftmodetrue" in the main .tex BEFORE % latex-rules — reusable preamble (see the latex-rules skill).
% Reusable macros only; define paper-specific notation in your main .tex, not here.

% ==== packages ====

% ---- fonts & encoding ----
\usepackage[T1]{fontenc}   % proper 8-bit fonts + real hyphenation
\usepackage{lmodern}       % scalable Latin Modern (fixes bitmap CM warnings)

% ---- math ----
\usepackage{amsmath}
\usepackage{amssymb}
\usepackage{mathtools}     % loads + extends amsmath
\usepackage{dsfont}        % \mathds{1}

% ---- graphics & drawing ----
\usepackage{graphicx}
\usepackage[export]{adjustbox}
\usepackage{tikz}
\usetikzlibrary{shapes}
\usetikzlibrary{positioning}
\usetikzlibrary{trees}
%\usetikzlibrary{graphs,graphdrawing}
%\usegdlibrary{trees}
\usepackage[most]{tcolorbox}
\usepackage{subcaption}    % NOT subfigure / subfig
\usepackage{pdflscape}

% ---- color ----
\usepackage{xcolor}

% ---- tables ----
\usepackage{booktabs}
\usepackage{multirow}
\usepackage{makecell}
\usepackage{colortbl}

% ---- typography, lists, units ----
\usepackage[protrusion=true,expansion=true]{microtype}  % char protrusion + font expansion; drop expansion under XeLaTeX (unsupported)
\usepackage{csquotes}
\usepackage{enumitem}
\usepackage{siunitx}
\sisetup{detect-all, per-mode=symbol}   % detect-all: match surrounding font; per-mode=symbol: m/s not m s^-1
\usepackage{xspace}

% ---- floats, algorithms, symbols ----
\usepackage{float}
\usepackage{afterpage}
\usepackage{footmisc}
\usepackage{listings}
\usepackage{algorithm}
\usepackage{algpseudocode}   % algorithmicx; modern replacement for algorithmic
\usepackage{pifont}
%\usepackage{pgfplots}

% ---- occasionally useful (uncomment as needed) ----
%\usepackage{tabu}
%\usepackage[defaultlines=3,all]{nowidow}
%\usepackage{setspace}
%\usepackage{geometry}
%\usepackage[all]{xy}
%\usepackage{natbib}
%\usepackage{subfigure}
% \usepackage{svg}   % needs --shell-escape + inkscape and is rejected by arXiv; pre-convert SVGs to PDF instead

% ==== math notation ====
% Symbol macros are wrapped in \ensuremath, so they also work in text (\R, \bx, ... auto-enter math).

% ---- quote marks in math ----
\DeclareMathSymbol{\mlq}{\mathord}{operators}{``}
\DeclareMathSymbol{\mrq}{\mathord}{operators}{`'}

% ---- number sets (blackboard bold) ----
% NOTE: \P overrides the built-in pilcrow (¶); see also \S \O \L \H below.
\newcommand{\R}{\ensuremath{\mathbb{R}}}
\newcommand{\N}{\ensuremath{\mathbb{N}}}
\newcommand{\Z}{\ensuremath{\mathbb{Z}}}
\newcommand{\Q}{\ensuremath{\mathbb{Q}}}
\newcommand{\C}{\ensuremath{\mathbb{C}}}
\newcommand{\D}{\ensuremath{\mathbb{D}}}
\newcommand{\E}{\ensuremath{\mathbb{E}}}
\newcommand{\X}{\ensuremath{\mathbb{X}}}
\newcommand{\Y}{\ensuremath{\mathbb{Y}}}
\renewcommand{\P}{\ensuremath{\mathbb{P}}}
\newcommand{\sigB}{\ensuremath{\mathbb{B}}}
\newcommand{\1}{\ensuremath{\mathds{1}}}

% ---- calligraphic sets ----
% NOTE: \S \O \L \H override built-in text commands (§, Ø, Ł, and the o-double-acute accent \H{o}).
\newcommand{\calN}{\ensuremath{\mathcal{N}}}
\newcommand{\calD}{\ensuremath{\mathcal{D}}}
\newcommand{\A}{\ensuremath{\mathcal{A}}}
\newcommand{\B}{\ensuremath{\mathcal{B}}}
\newcommand{\I}{\ensuremath{\mathcal{I}}}
\newcommand{\F}{\ensuremath{\mathcal{F}}}
\newcommand{\G}{\ensuremath{\mathcal{G}}}
\newcommand{\M}{\ensuremath{\mathcal{M}}}
\newcommand{\Tau}{\ensuremath{\mathcal{T}}}
\newcommand{\W}{\ensuremath{\mathcal{W}}}
\newcommand{\sigE}{\ensuremath{\mathcal{E}}}
\renewcommand{\H}{\ensuremath{\mathcal{H}}}
\renewcommand{\L}{\ensuremath{\mathcal{L}}}
\renewcommand{\S}{\ensuremath{\mathcal{S}}}
\renewcommand{\O}{\ensuremath{\mathcal{O}}}

% ---- named operators ----
% NOTE: \Re \Im are redefined to roman "Re"/"Im" (not the default blackletter ℜ/ℑ).
\DeclareMathOperator{\softmax}{softmax}
\DeclareMathOperator*{\argmax}{arg\,max}   % * = subscript sits underneath in display
\DeclareMathOperator{\Var}{Var}
\newcommand{\cov}{\ensuremath{\operatorname{cov}}}
\newcommand{\corr}{\ensuremath{\operatorname{corr}}}
\newcommand{\spanop}{\ensuremath{\operatorname{span}}}
\newcommand{\att}{\ensuremath{\operatorname{attention}}}
\newcommand{\suff}{\ensuremath{\operatorname{suff}}}
\newcommand{\comp}{\ensuremath{\operatorname{comp}}}
\newcommand{\In}{\ensuremath{\operatorname{In}}}
\newcommand{\opcap}{\ensuremath{\operatorname{cap}}}
\newcommand{\entr}{\ensuremath{\operatorname{entries}}}
\newcommand{\ops}{\ensuremath{\operatorname{ops}}}
\newcommand{\id}{\ensuremath{\mathrm{id}}}
\renewcommand{\Re}{\ensuremath{\operatorname{Re}}}
\renewcommand{\Im}{\ensuremath{\operatorname{Im}}}

% ---- symbols & shorthands ----
\newcommand{\eps}{\ensuremath{\varepsilon}}
\renewcommand{\phi}{\ensuremath{\varphi}}     % deliberate: \phi -> varphi
\newcommand{\del}{\ensuremath{\partial}}
\newcommand{\grad}{\ensuremath{\nabla}}
\newcommand{\indep}{\ensuremath{\perp\!\!\!\perp}}
\newcommand{\emptyword}{\ensuremath{\varnothing}}
\newcommand{\tensor}{\ensuremath{\otimes}}
\newcommand{\bigtensor}{\ensuremath{\bigotimes}}
\newcommand{\To}{\ensuremath{\longrightarrow}}
\newcommand{\eqover}[1]{\ensuremath{\stackrel{#1}{=}}}
\newcommand{\darover}[1]{\ensuremath{\xrightarrow{#1}}}
\newcommand{\Hoel}{\text{Höl}}

% ---- paired delimiters ----
\DeclarePairedDelimiter{\norm}{\lVert}{\rVert}   % \norm{x}, \norm*{x} (full size), \norm[\big]{x}
\DeclarePairedDelimiter{\abs}{\lvert}{\rvert}

% ---- bold vectors / matrices ----
\newcommand{\bx}{\ensuremath{\mathbf{x}}}
\newcommand{\bX}{\ensuremath{\mathbf{X}}}
\newcommand{\by}{\ensuremath{\mathbf{y}}}
\newcommand{\bz}{\ensuremath{\mathbf{z}}}
\newcommand{\bB}{\ensuremath{\mathbf{B}}}
\newcommand{\bp}{\ensuremath{\mathbf{p}}}
\newcommand{\bW}{\ensuremath{\mathbf{W}}}
\newcommand{\bt}{\ensuremath{\mathbf{t}}}
\newcommand{\be}{\ensuremath{\mathbf{e}}}

% ---- small constructs ----
\newcommand{\Ball}[2]{\ensuremath{B_#1(#2)}}   % \Ball{\eps}{center}
\newcommand{\twobytwo}[4]{\ensuremath{\begin{pmatrix} #1 & #2 \\ #3 & #4 \end{pmatrix}}}

% ==== text, color & layout helpers ====
\definecolor{darkgreen}{HTML}{009B55}
\newcommand{\gtxt}[1]{\text{\textcolor{gray}{#1}}}
\newcommand{\grntxt}[1]{\text{\textcolor{darkgreen}{#1}}}
\newcommand{\rdtxt}[1]{\text{\textcolor{red}{#1}}}
\newcommand{\code}[1]{\texttt{#1}}
\newcommand{\tldr}{\textbf{TL;DR:}\xspace}
\newcommand{\parahead}[1]{\noindent\textbf{#1.}\enspace}  % run-in heading; use instead of \paragraph to save space
\newcommand*\rot{\rotatebox{90}}

% ---- table cell helpers ----
\newcommand{\cmark}{\ding{51}}
\newcommand{\xmark}{\ding{55}}
\newcommand{\rowhighlight}{\rowcolor{gray!30}}
\newcommand{\cellhighlight}{\cellcolor{yellow!40}}

% ---- rebuttal reviewer tags: bold, distinct colors to signpost which reviewer a point answers ----
\newcommand{\Rone}{\textbf{\textcolor[HTML]{1F77B4}{R1}}}    % blue
\newcommand{\Rtwo}{\textbf{\textcolor[HTML]{D62728}{R2}}}    % red
\newcommand{\Rthree}{\textbf{\textcolor[HTML]{2CA02C}{R3}}}  % green
\newcommand{\Rfour}{\textbf{\textcolor[HTML]{D9730D}{R4}}}   % muted orange

% ---- takeaway callout box ----
\newtcolorbox{takeaway}[1][Takeaway]{
  colback=blue!4!white,      % very light fill
  colframe=blue!50!black,    % muted dark-blue frame
  boxrule=0.6pt,
  arc=2pt,
  left=6pt, right=6pt, top=2pt, bottom=2pt,
  title=#1,
  fonttitle=\bfseries\small,
  coltitle=white,
  colbacktitle=blue!50!black
}

% ==== links & smart references: load LAST (delete if your class already loads hyperref) ====
\usepackage[hidelinks]{hyperref}
\usepackage[capitalize,noabbrev]{cleveref}   % \cref{...}, \Cref{...} — must come after hyperref

% ==== draft mode & \todo ====
% "Explicit" = the main .tex defined \ifdraftmode before \input{packages}.
% (The \ifcsname/\csname dance avoids skipping over an already-defined \ifdraftmode, which
%  crashes the naive \ifdefined...\newif idiom.)
\newif\ifdraftmodeexplicit
\ifcsname ifdraftmode\endcsname \draftmodeexplicittrue \else \draftmodeexplicitfalse \fi
\ifdraftmodeexplicit\else \expandafter\newif\csname ifdraftmode\endcsname \fi

% \todo: visible red box by default and in draft mode; a no-op only when draft mode is
% explicitly OFF (\draftmodefalse in the main .tex), so leftover todos vanish in the final build.
\newcommand{\todo}[1]{\colorbox{red}{TODO: #1}}
\ifdraftmodeexplicit \ifdraftmode\else \renewcommand{\todo}[1]{}\fi \fi

% Draft-mode aids (line numbers, label keys, overfull bars) — on only when draft mode is enabled.
% Worth it for long documents / theses; usually overkill for a short single paper.
% Enable by putting "\newif\ifdraftmode \draftmodetrue" in the main .tex BEFORE \input{packages}.
\ifdraftmode
  \usepackage{lineno}\linenumbers   % line numbers for "see line 214" review comments (prose only; align/equation bodies aren't numbered without extra setup)
  \usepackage{showkeys}             % print \label / \ref keys in the margin
  \overfullrule=5pt                 % black bar in the margin marks every overfull line
\fi
.
\ifdraftmode
  \usepackage{lineno}\linenumbers   % line numbers for "see line 214" review comments (prose only; align/equation bodies aren't numbered without extra setup)
  \usepackage{showkeys}             % print \label / \ref keys in the margin
  \overfullrule=5pt                 % black bar in the margin marks every overfull line
\fi
.
\ifdraftmode
  \usepackage{lineno}\linenumbers   % line numbers for "see line 214" review comments (prose only; align/equation bodies aren't numbered without extra setup)
  \usepackage{showkeys}             % print \label / \ref keys in the margin
  \overfullrule=5pt                 % black bar in the margin marks every overfull line
\fi
.
% (The \ifcsname/\csname dance avoids skipping over an already-defined \ifdraftmode, which
%  crashes the naive \ifdefined...\newif idiom.)
\newif\ifdraftmodeexplicit
\ifcsname ifdraftmode\endcsname \draftmodeexplicittrue \else \draftmodeexplicitfalse \fi
\ifdraftmodeexplicit\else \expandafter\newif\csname ifdraftmode\endcsname \fi

% \todo: visible red box by default and in draft mode; a no-op only when draft mode is
% explicitly OFF (\draftmodefalse in the main .tex), so leftover todos vanish in the final build.
\newcommand{\todo}[1]{\colorbox{red}{TODO: #1}}
\ifdraftmodeexplicit \ifdraftmode\else \renewcommand{\todo}[1]{}\fi \fi

% Draft-mode aids (line numbers, label keys, overfull bars) — on only when draft mode is enabled.
% Worth it for long documents / theses; usually overkill for a short single paper.
% Enable by putting "\newif\ifdraftmode \draftmodetrue" in the main .tex BEFORE % latex-rules — reusable preamble (see the latex-rules skill).
% Reusable macros only; define paper-specific notation in your main .tex, not here.

% ==== packages ====

% ---- fonts & encoding ----
\usepackage[T1]{fontenc}   % proper 8-bit fonts + real hyphenation
\usepackage{lmodern}       % scalable Latin Modern (fixes bitmap CM warnings)

% ---- math ----
\usepackage{amsmath}
\usepackage{amssymb}
\usepackage{mathtools}     % loads + extends amsmath
\usepackage{dsfont}        % \mathds{1}

% ---- graphics & drawing ----
\usepackage{graphicx}
\usepackage[export]{adjustbox}
\usepackage{tikz}
\usetikzlibrary{shapes}
\usetikzlibrary{positioning}
\usetikzlibrary{trees}
%\usetikzlibrary{graphs,graphdrawing}
%\usegdlibrary{trees}
\usepackage[most]{tcolorbox}
\usepackage{subcaption}    % NOT subfigure / subfig
\usepackage{pdflscape}

% ---- color ----
\usepackage{xcolor}

% ---- tables ----
\usepackage{booktabs}
\usepackage{multirow}
\usepackage{makecell}
\usepackage{colortbl}

% ---- typography, lists, units ----
\usepackage[protrusion=true,expansion=true]{microtype}  % char protrusion + font expansion; drop expansion under XeLaTeX (unsupported)
\usepackage{csquotes}
\usepackage{enumitem}
\usepackage{siunitx}
\sisetup{detect-all, per-mode=symbol}   % detect-all: match surrounding font; per-mode=symbol: m/s not m s^-1
\usepackage{xspace}

% ---- floats, algorithms, symbols ----
\usepackage{float}
\usepackage{afterpage}
\usepackage{footmisc}
\usepackage{listings}
\usepackage{algorithm}
\usepackage{algpseudocode}   % algorithmicx; modern replacement for algorithmic
\usepackage{pifont}
%\usepackage{pgfplots}

% ---- occasionally useful (uncomment as needed) ----
%\usepackage{tabu}
%\usepackage[defaultlines=3,all]{nowidow}
%\usepackage{setspace}
%\usepackage{geometry}
%\usepackage[all]{xy}
%\usepackage{natbib}
%\usepackage{subfigure}
% \usepackage{svg}   % needs --shell-escape + inkscape and is rejected by arXiv; pre-convert SVGs to PDF instead

% ==== math notation ====
% Symbol macros are wrapped in \ensuremath, so they also work in text (\R, \bx, ... auto-enter math).

% ---- quote marks in math ----
\DeclareMathSymbol{\mlq}{\mathord}{operators}{``}
\DeclareMathSymbol{\mrq}{\mathord}{operators}{`'}

% ---- number sets (blackboard bold) ----
% NOTE: \P overrides the built-in pilcrow (¶); see also \S \O \L \H below.
\newcommand{\R}{\ensuremath{\mathbb{R}}}
\newcommand{\N}{\ensuremath{\mathbb{N}}}
\newcommand{\Z}{\ensuremath{\mathbb{Z}}}
\newcommand{\Q}{\ensuremath{\mathbb{Q}}}
\newcommand{\C}{\ensuremath{\mathbb{C}}}
\newcommand{\D}{\ensuremath{\mathbb{D}}}
\newcommand{\E}{\ensuremath{\mathbb{E}}}
\newcommand{\X}{\ensuremath{\mathbb{X}}}
\newcommand{\Y}{\ensuremath{\mathbb{Y}}}
\renewcommand{\P}{\ensuremath{\mathbb{P}}}
\newcommand{\sigB}{\ensuremath{\mathbb{B}}}
\newcommand{\1}{\ensuremath{\mathds{1}}}

% ---- calligraphic sets ----
% NOTE: \S \O \L \H override built-in text commands (§, Ø, Ł, and the o-double-acute accent \H{o}).
\newcommand{\calN}{\ensuremath{\mathcal{N}}}
\newcommand{\calD}{\ensuremath{\mathcal{D}}}
\newcommand{\A}{\ensuremath{\mathcal{A}}}
\newcommand{\B}{\ensuremath{\mathcal{B}}}
\newcommand{\I}{\ensuremath{\mathcal{I}}}
\newcommand{\F}{\ensuremath{\mathcal{F}}}
\newcommand{\G}{\ensuremath{\mathcal{G}}}
\newcommand{\M}{\ensuremath{\mathcal{M}}}
\newcommand{\Tau}{\ensuremath{\mathcal{T}}}
\newcommand{\W}{\ensuremath{\mathcal{W}}}
\newcommand{\sigE}{\ensuremath{\mathcal{E}}}
\renewcommand{\H}{\ensuremath{\mathcal{H}}}
\renewcommand{\L}{\ensuremath{\mathcal{L}}}
\renewcommand{\S}{\ensuremath{\mathcal{S}}}
\renewcommand{\O}{\ensuremath{\mathcal{O}}}

% ---- named operators ----
% NOTE: \Re \Im are redefined to roman "Re"/"Im" (not the default blackletter ℜ/ℑ).
\DeclareMathOperator{\softmax}{softmax}
\DeclareMathOperator*{\argmax}{arg\,max}   % * = subscript sits underneath in display
\DeclareMathOperator{\Var}{Var}
\newcommand{\cov}{\ensuremath{\operatorname{cov}}}
\newcommand{\corr}{\ensuremath{\operatorname{corr}}}
\newcommand{\spanop}{\ensuremath{\operatorname{span}}}
\newcommand{\att}{\ensuremath{\operatorname{attention}}}
\newcommand{\suff}{\ensuremath{\operatorname{suff}}}
\newcommand{\comp}{\ensuremath{\operatorname{comp}}}
\newcommand{\In}{\ensuremath{\operatorname{In}}}
\newcommand{\opcap}{\ensuremath{\operatorname{cap}}}
\newcommand{\entr}{\ensuremath{\operatorname{entries}}}
\newcommand{\ops}{\ensuremath{\operatorname{ops}}}
\newcommand{\id}{\ensuremath{\mathrm{id}}}
\renewcommand{\Re}{\ensuremath{\operatorname{Re}}}
\renewcommand{\Im}{\ensuremath{\operatorname{Im}}}

% ---- symbols & shorthands ----
\newcommand{\eps}{\ensuremath{\varepsilon}}
\renewcommand{\phi}{\ensuremath{\varphi}}     % deliberate: \phi -> varphi
\newcommand{\del}{\ensuremath{\partial}}
\newcommand{\grad}{\ensuremath{\nabla}}
\newcommand{\indep}{\ensuremath{\perp\!\!\!\perp}}
\newcommand{\emptyword}{\ensuremath{\varnothing}}
\newcommand{\tensor}{\ensuremath{\otimes}}
\newcommand{\bigtensor}{\ensuremath{\bigotimes}}
\newcommand{\To}{\ensuremath{\longrightarrow}}
\newcommand{\eqover}[1]{\ensuremath{\stackrel{#1}{=}}}
\newcommand{\darover}[1]{\ensuremath{\xrightarrow{#1}}}
\newcommand{\Hoel}{\text{Höl}}

% ---- paired delimiters ----
\DeclarePairedDelimiter{\norm}{\lVert}{\rVert}   % \norm{x}, \norm*{x} (full size), \norm[\big]{x}
\DeclarePairedDelimiter{\abs}{\lvert}{\rvert}

% ---- bold vectors / matrices ----
\newcommand{\bx}{\ensuremath{\mathbf{x}}}
\newcommand{\bX}{\ensuremath{\mathbf{X}}}
\newcommand{\by}{\ensuremath{\mathbf{y}}}
\newcommand{\bz}{\ensuremath{\mathbf{z}}}
\newcommand{\bB}{\ensuremath{\mathbf{B}}}
\newcommand{\bp}{\ensuremath{\mathbf{p}}}
\newcommand{\bW}{\ensuremath{\mathbf{W}}}
\newcommand{\bt}{\ensuremath{\mathbf{t}}}
\newcommand{\be}{\ensuremath{\mathbf{e}}}

% ---- small constructs ----
\newcommand{\Ball}[2]{\ensuremath{B_#1(#2)}}   % \Ball{\eps}{center}
\newcommand{\twobytwo}[4]{\ensuremath{\begin{pmatrix} #1 & #2 \\ #3 & #4 \end{pmatrix}}}

% ==== text, color & layout helpers ====
\definecolor{darkgreen}{HTML}{009B55}
\newcommand{\gtxt}[1]{\text{\textcolor{gray}{#1}}}
\newcommand{\grntxt}[1]{\text{\textcolor{darkgreen}{#1}}}
\newcommand{\rdtxt}[1]{\text{\textcolor{red}{#1}}}
\newcommand{\code}[1]{\texttt{#1}}
\newcommand{\tldr}{\textbf{TL;DR:}\xspace}
\newcommand{\parahead}[1]{\noindent\textbf{#1.}\enspace}  % run-in heading; use instead of \paragraph to save space
\newcommand*\rot{\rotatebox{90}}

% ---- table cell helpers ----
\newcommand{\cmark}{\ding{51}}
\newcommand{\xmark}{\ding{55}}
\newcommand{\rowhighlight}{\rowcolor{gray!30}}
\newcommand{\cellhighlight}{\cellcolor{yellow!40}}

% ---- rebuttal reviewer tags: bold, distinct colors to signpost which reviewer a point answers ----
\newcommand{\Rone}{\textbf{\textcolor[HTML]{1F77B4}{R1}}}    % blue
\newcommand{\Rtwo}{\textbf{\textcolor[HTML]{D62728}{R2}}}    % red
\newcommand{\Rthree}{\textbf{\textcolor[HTML]{2CA02C}{R3}}}  % green
\newcommand{\Rfour}{\textbf{\textcolor[HTML]{D9730D}{R4}}}   % muted orange

% ---- takeaway callout box ----
\newtcolorbox{takeaway}[1][Takeaway]{
  colback=blue!4!white,      % very light fill
  colframe=blue!50!black,    % muted dark-blue frame
  boxrule=0.6pt,
  arc=2pt,
  left=6pt, right=6pt, top=2pt, bottom=2pt,
  title=#1,
  fonttitle=\bfseries\small,
  coltitle=white,
  colbacktitle=blue!50!black
}

% ==== links & smart references: load LAST (delete if your class already loads hyperref) ====
\usepackage[hidelinks]{hyperref}
\usepackage[capitalize,noabbrev]{cleveref}   % \cref{...}, \Cref{...} — must come after hyperref

% ==== draft mode & \todo ====
% "Explicit" = the main .tex defined \ifdraftmode before % latex-rules — reusable preamble (see the latex-rules skill).
% Reusable macros only; define paper-specific notation in your main .tex, not here.

% ==== packages ====

% ---- fonts & encoding ----
\usepackage[T1]{fontenc}   % proper 8-bit fonts + real hyphenation
\usepackage{lmodern}       % scalable Latin Modern (fixes bitmap CM warnings)

% ---- math ----
\usepackage{amsmath}
\usepackage{amssymb}
\usepackage{mathtools}     % loads + extends amsmath
\usepackage{dsfont}        % \mathds{1}

% ---- graphics & drawing ----
\usepackage{graphicx}
\usepackage[export]{adjustbox}
\usepackage{tikz}
\usetikzlibrary{shapes}
\usetikzlibrary{positioning}
\usetikzlibrary{trees}
%\usetikzlibrary{graphs,graphdrawing}
%\usegdlibrary{trees}
\usepackage[most]{tcolorbox}
\usepackage{subcaption}    % NOT subfigure / subfig
\usepackage{pdflscape}

% ---- color ----
\usepackage{xcolor}

% ---- tables ----
\usepackage{booktabs}
\usepackage{multirow}
\usepackage{makecell}
\usepackage{colortbl}

% ---- typography, lists, units ----
\usepackage[protrusion=true,expansion=true]{microtype}  % char protrusion + font expansion; drop expansion under XeLaTeX (unsupported)
\usepackage{csquotes}
\usepackage{enumitem}
\usepackage{siunitx}
\sisetup{detect-all, per-mode=symbol}   % detect-all: match surrounding font; per-mode=symbol: m/s not m s^-1
\usepackage{xspace}

% ---- floats, algorithms, symbols ----
\usepackage{float}
\usepackage{afterpage}
\usepackage{footmisc}
\usepackage{listings}
\usepackage{algorithm}
\usepackage{algpseudocode}   % algorithmicx; modern replacement for algorithmic
\usepackage{pifont}
%\usepackage{pgfplots}

% ---- occasionally useful (uncomment as needed) ----
%\usepackage{tabu}
%\usepackage[defaultlines=3,all]{nowidow}
%\usepackage{setspace}
%\usepackage{geometry}
%\usepackage[all]{xy}
%\usepackage{natbib}
%\usepackage{subfigure}
% \usepackage{svg}   % needs --shell-escape + inkscape and is rejected by arXiv; pre-convert SVGs to PDF instead

% ==== math notation ====
% Symbol macros are wrapped in \ensuremath, so they also work in text (\R, \bx, ... auto-enter math).

% ---- quote marks in math ----
\DeclareMathSymbol{\mlq}{\mathord}{operators}{``}
\DeclareMathSymbol{\mrq}{\mathord}{operators}{`'}

% ---- number sets (blackboard bold) ----
% NOTE: \P overrides the built-in pilcrow (¶); see also \S \O \L \H below.
\newcommand{\R}{\ensuremath{\mathbb{R}}}
\newcommand{\N}{\ensuremath{\mathbb{N}}}
\newcommand{\Z}{\ensuremath{\mathbb{Z}}}
\newcommand{\Q}{\ensuremath{\mathbb{Q}}}
\newcommand{\C}{\ensuremath{\mathbb{C}}}
\newcommand{\D}{\ensuremath{\mathbb{D}}}
\newcommand{\E}{\ensuremath{\mathbb{E}}}
\newcommand{\X}{\ensuremath{\mathbb{X}}}
\newcommand{\Y}{\ensuremath{\mathbb{Y}}}
\renewcommand{\P}{\ensuremath{\mathbb{P}}}
\newcommand{\sigB}{\ensuremath{\mathbb{B}}}
\newcommand{\1}{\ensuremath{\mathds{1}}}

% ---- calligraphic sets ----
% NOTE: \S \O \L \H override built-in text commands (§, Ø, Ł, and the o-double-acute accent \H{o}).
\newcommand{\calN}{\ensuremath{\mathcal{N}}}
\newcommand{\calD}{\ensuremath{\mathcal{D}}}
\newcommand{\A}{\ensuremath{\mathcal{A}}}
\newcommand{\B}{\ensuremath{\mathcal{B}}}
\newcommand{\I}{\ensuremath{\mathcal{I}}}
\newcommand{\F}{\ensuremath{\mathcal{F}}}
\newcommand{\G}{\ensuremath{\mathcal{G}}}
\newcommand{\M}{\ensuremath{\mathcal{M}}}
\newcommand{\Tau}{\ensuremath{\mathcal{T}}}
\newcommand{\W}{\ensuremath{\mathcal{W}}}
\newcommand{\sigE}{\ensuremath{\mathcal{E}}}
\renewcommand{\H}{\ensuremath{\mathcal{H}}}
\renewcommand{\L}{\ensuremath{\mathcal{L}}}
\renewcommand{\S}{\ensuremath{\mathcal{S}}}
\renewcommand{\O}{\ensuremath{\mathcal{O}}}

% ---- named operators ----
% NOTE: \Re \Im are redefined to roman "Re"/"Im" (not the default blackletter ℜ/ℑ).
\DeclareMathOperator{\softmax}{softmax}
\DeclareMathOperator*{\argmax}{arg\,max}   % * = subscript sits underneath in display
\DeclareMathOperator{\Var}{Var}
\newcommand{\cov}{\ensuremath{\operatorname{cov}}}
\newcommand{\corr}{\ensuremath{\operatorname{corr}}}
\newcommand{\spanop}{\ensuremath{\operatorname{span}}}
\newcommand{\att}{\ensuremath{\operatorname{attention}}}
\newcommand{\suff}{\ensuremath{\operatorname{suff}}}
\newcommand{\comp}{\ensuremath{\operatorname{comp}}}
\newcommand{\In}{\ensuremath{\operatorname{In}}}
\newcommand{\opcap}{\ensuremath{\operatorname{cap}}}
\newcommand{\entr}{\ensuremath{\operatorname{entries}}}
\newcommand{\ops}{\ensuremath{\operatorname{ops}}}
\newcommand{\id}{\ensuremath{\mathrm{id}}}
\renewcommand{\Re}{\ensuremath{\operatorname{Re}}}
\renewcommand{\Im}{\ensuremath{\operatorname{Im}}}

% ---- symbols & shorthands ----
\newcommand{\eps}{\ensuremath{\varepsilon}}
\renewcommand{\phi}{\ensuremath{\varphi}}     % deliberate: \phi -> varphi
\newcommand{\del}{\ensuremath{\partial}}
\newcommand{\grad}{\ensuremath{\nabla}}
\newcommand{\indep}{\ensuremath{\perp\!\!\!\perp}}
\newcommand{\emptyword}{\ensuremath{\varnothing}}
\newcommand{\tensor}{\ensuremath{\otimes}}
\newcommand{\bigtensor}{\ensuremath{\bigotimes}}
\newcommand{\To}{\ensuremath{\longrightarrow}}
\newcommand{\eqover}[1]{\ensuremath{\stackrel{#1}{=}}}
\newcommand{\darover}[1]{\ensuremath{\xrightarrow{#1}}}
\newcommand{\Hoel}{\text{Höl}}

% ---- paired delimiters ----
\DeclarePairedDelimiter{\norm}{\lVert}{\rVert}   % \norm{x}, \norm*{x} (full size), \norm[\big]{x}
\DeclarePairedDelimiter{\abs}{\lvert}{\rvert}

% ---- bold vectors / matrices ----
\newcommand{\bx}{\ensuremath{\mathbf{x}}}
\newcommand{\bX}{\ensuremath{\mathbf{X}}}
\newcommand{\by}{\ensuremath{\mathbf{y}}}
\newcommand{\bz}{\ensuremath{\mathbf{z}}}
\newcommand{\bB}{\ensuremath{\mathbf{B}}}
\newcommand{\bp}{\ensuremath{\mathbf{p}}}
\newcommand{\bW}{\ensuremath{\mathbf{W}}}
\newcommand{\bt}{\ensuremath{\mathbf{t}}}
\newcommand{\be}{\ensuremath{\mathbf{e}}}

% ---- small constructs ----
\newcommand{\Ball}[2]{\ensuremath{B_#1(#2)}}   % \Ball{\eps}{center}
\newcommand{\twobytwo}[4]{\ensuremath{\begin{pmatrix} #1 & #2 \\ #3 & #4 \end{pmatrix}}}

% ==== text, color & layout helpers ====
\definecolor{darkgreen}{HTML}{009B55}
\newcommand{\gtxt}[1]{\text{\textcolor{gray}{#1}}}
\newcommand{\grntxt}[1]{\text{\textcolor{darkgreen}{#1}}}
\newcommand{\rdtxt}[1]{\text{\textcolor{red}{#1}}}
\newcommand{\code}[1]{\texttt{#1}}
\newcommand{\tldr}{\textbf{TL;DR:}\xspace}
\newcommand{\parahead}[1]{\noindent\textbf{#1.}\enspace}  % run-in heading; use instead of \paragraph to save space
\newcommand*\rot{\rotatebox{90}}

% ---- table cell helpers ----
\newcommand{\cmark}{\ding{51}}
\newcommand{\xmark}{\ding{55}}
\newcommand{\rowhighlight}{\rowcolor{gray!30}}
\newcommand{\cellhighlight}{\cellcolor{yellow!40}}

% ---- rebuttal reviewer tags: bold, distinct colors to signpost which reviewer a point answers ----
\newcommand{\Rone}{\textbf{\textcolor[HTML]{1F77B4}{R1}}}    % blue
\newcommand{\Rtwo}{\textbf{\textcolor[HTML]{D62728}{R2}}}    % red
\newcommand{\Rthree}{\textbf{\textcolor[HTML]{2CA02C}{R3}}}  % green
\newcommand{\Rfour}{\textbf{\textcolor[HTML]{D9730D}{R4}}}   % muted orange

% ---- takeaway callout box ----
\newtcolorbox{takeaway}[1][Takeaway]{
  colback=blue!4!white,      % very light fill
  colframe=blue!50!black,    % muted dark-blue frame
  boxrule=0.6pt,
  arc=2pt,
  left=6pt, right=6pt, top=2pt, bottom=2pt,
  title=#1,
  fonttitle=\bfseries\small,
  coltitle=white,
  colbacktitle=blue!50!black
}

% ==== links & smart references: load LAST (delete if your class already loads hyperref) ====
\usepackage[hidelinks]{hyperref}
\usepackage[capitalize,noabbrev]{cleveref}   % \cref{...}, \Cref{...} — must come after hyperref

% ==== draft mode & \todo ====
% "Explicit" = the main .tex defined \ifdraftmode before % latex-rules — reusable preamble (see the latex-rules skill).
% Reusable macros only; define paper-specific notation in your main .tex, not here.

% ==== packages ====

% ---- fonts & encoding ----
\usepackage[T1]{fontenc}   % proper 8-bit fonts + real hyphenation
\usepackage{lmodern}       % scalable Latin Modern (fixes bitmap CM warnings)

% ---- math ----
\usepackage{amsmath}
\usepackage{amssymb}
\usepackage{mathtools}     % loads + extends amsmath
\usepackage{dsfont}        % \mathds{1}

% ---- graphics & drawing ----
\usepackage{graphicx}
\usepackage[export]{adjustbox}
\usepackage{tikz}
\usetikzlibrary{shapes}
\usetikzlibrary{positioning}
\usetikzlibrary{trees}
%\usetikzlibrary{graphs,graphdrawing}
%\usegdlibrary{trees}
\usepackage[most]{tcolorbox}
\usepackage{subcaption}    % NOT subfigure / subfig
\usepackage{pdflscape}

% ---- color ----
\usepackage{xcolor}

% ---- tables ----
\usepackage{booktabs}
\usepackage{multirow}
\usepackage{makecell}
\usepackage{colortbl}

% ---- typography, lists, units ----
\usepackage[protrusion=true,expansion=true]{microtype}  % char protrusion + font expansion; drop expansion under XeLaTeX (unsupported)
\usepackage{csquotes}
\usepackage{enumitem}
\usepackage{siunitx}
\sisetup{detect-all, per-mode=symbol}   % detect-all: match surrounding font; per-mode=symbol: m/s not m s^-1
\usepackage{xspace}

% ---- floats, algorithms, symbols ----
\usepackage{float}
\usepackage{afterpage}
\usepackage{footmisc}
\usepackage{listings}
\usepackage{algorithm}
\usepackage{algpseudocode}   % algorithmicx; modern replacement for algorithmic
\usepackage{pifont}
%\usepackage{pgfplots}

% ---- occasionally useful (uncomment as needed) ----
%\usepackage{tabu}
%\usepackage[defaultlines=3,all]{nowidow}
%\usepackage{setspace}
%\usepackage{geometry}
%\usepackage[all]{xy}
%\usepackage{natbib}
%\usepackage{subfigure}
% \usepackage{svg}   % needs --shell-escape + inkscape and is rejected by arXiv; pre-convert SVGs to PDF instead

% ==== math notation ====
% Symbol macros are wrapped in \ensuremath, so they also work in text (\R, \bx, ... auto-enter math).

% ---- quote marks in math ----
\DeclareMathSymbol{\mlq}{\mathord}{operators}{``}
\DeclareMathSymbol{\mrq}{\mathord}{operators}{`'}

% ---- number sets (blackboard bold) ----
% NOTE: \P overrides the built-in pilcrow (¶); see also \S \O \L \H below.
\newcommand{\R}{\ensuremath{\mathbb{R}}}
\newcommand{\N}{\ensuremath{\mathbb{N}}}
\newcommand{\Z}{\ensuremath{\mathbb{Z}}}
\newcommand{\Q}{\ensuremath{\mathbb{Q}}}
\newcommand{\C}{\ensuremath{\mathbb{C}}}
\newcommand{\D}{\ensuremath{\mathbb{D}}}
\newcommand{\E}{\ensuremath{\mathbb{E}}}
\newcommand{\X}{\ensuremath{\mathbb{X}}}
\newcommand{\Y}{\ensuremath{\mathbb{Y}}}
\renewcommand{\P}{\ensuremath{\mathbb{P}}}
\newcommand{\sigB}{\ensuremath{\mathbb{B}}}
\newcommand{\1}{\ensuremath{\mathds{1}}}

% ---- calligraphic sets ----
% NOTE: \S \O \L \H override built-in text commands (§, Ø, Ł, and the o-double-acute accent \H{o}).
\newcommand{\calN}{\ensuremath{\mathcal{N}}}
\newcommand{\calD}{\ensuremath{\mathcal{D}}}
\newcommand{\A}{\ensuremath{\mathcal{A}}}
\newcommand{\B}{\ensuremath{\mathcal{B}}}
\newcommand{\I}{\ensuremath{\mathcal{I}}}
\newcommand{\F}{\ensuremath{\mathcal{F}}}
\newcommand{\G}{\ensuremath{\mathcal{G}}}
\newcommand{\M}{\ensuremath{\mathcal{M}}}
\newcommand{\Tau}{\ensuremath{\mathcal{T}}}
\newcommand{\W}{\ensuremath{\mathcal{W}}}
\newcommand{\sigE}{\ensuremath{\mathcal{E}}}
\renewcommand{\H}{\ensuremath{\mathcal{H}}}
\renewcommand{\L}{\ensuremath{\mathcal{L}}}
\renewcommand{\S}{\ensuremath{\mathcal{S}}}
\renewcommand{\O}{\ensuremath{\mathcal{O}}}

% ---- named operators ----
% NOTE: \Re \Im are redefined to roman "Re"/"Im" (not the default blackletter ℜ/ℑ).
\DeclareMathOperator{\softmax}{softmax}
\DeclareMathOperator*{\argmax}{arg\,max}   % * = subscript sits underneath in display
\DeclareMathOperator{\Var}{Var}
\newcommand{\cov}{\ensuremath{\operatorname{cov}}}
\newcommand{\corr}{\ensuremath{\operatorname{corr}}}
\newcommand{\spanop}{\ensuremath{\operatorname{span}}}
\newcommand{\att}{\ensuremath{\operatorname{attention}}}
\newcommand{\suff}{\ensuremath{\operatorname{suff}}}
\newcommand{\comp}{\ensuremath{\operatorname{comp}}}
\newcommand{\In}{\ensuremath{\operatorname{In}}}
\newcommand{\opcap}{\ensuremath{\operatorname{cap}}}
\newcommand{\entr}{\ensuremath{\operatorname{entries}}}
\newcommand{\ops}{\ensuremath{\operatorname{ops}}}
\newcommand{\id}{\ensuremath{\mathrm{id}}}
\renewcommand{\Re}{\ensuremath{\operatorname{Re}}}
\renewcommand{\Im}{\ensuremath{\operatorname{Im}}}

% ---- symbols & shorthands ----
\newcommand{\eps}{\ensuremath{\varepsilon}}
\renewcommand{\phi}{\ensuremath{\varphi}}     % deliberate: \phi -> varphi
\newcommand{\del}{\ensuremath{\partial}}
\newcommand{\grad}{\ensuremath{\nabla}}
\newcommand{\indep}{\ensuremath{\perp\!\!\!\perp}}
\newcommand{\emptyword}{\ensuremath{\varnothing}}
\newcommand{\tensor}{\ensuremath{\otimes}}
\newcommand{\bigtensor}{\ensuremath{\bigotimes}}
\newcommand{\To}{\ensuremath{\longrightarrow}}
\newcommand{\eqover}[1]{\ensuremath{\stackrel{#1}{=}}}
\newcommand{\darover}[1]{\ensuremath{\xrightarrow{#1}}}
\newcommand{\Hoel}{\text{Höl}}

% ---- paired delimiters ----
\DeclarePairedDelimiter{\norm}{\lVert}{\rVert}   % \norm{x}, \norm*{x} (full size), \norm[\big]{x}
\DeclarePairedDelimiter{\abs}{\lvert}{\rvert}

% ---- bold vectors / matrices ----
\newcommand{\bx}{\ensuremath{\mathbf{x}}}
\newcommand{\bX}{\ensuremath{\mathbf{X}}}
\newcommand{\by}{\ensuremath{\mathbf{y}}}
\newcommand{\bz}{\ensuremath{\mathbf{z}}}
\newcommand{\bB}{\ensuremath{\mathbf{B}}}
\newcommand{\bp}{\ensuremath{\mathbf{p}}}
\newcommand{\bW}{\ensuremath{\mathbf{W}}}
\newcommand{\bt}{\ensuremath{\mathbf{t}}}
\newcommand{\be}{\ensuremath{\mathbf{e}}}

% ---- small constructs ----
\newcommand{\Ball}[2]{\ensuremath{B_#1(#2)}}   % \Ball{\eps}{center}
\newcommand{\twobytwo}[4]{\ensuremath{\begin{pmatrix} #1 & #2 \\ #3 & #4 \end{pmatrix}}}

% ==== text, color & layout helpers ====
\definecolor{darkgreen}{HTML}{009B55}
\newcommand{\gtxt}[1]{\text{\textcolor{gray}{#1}}}
\newcommand{\grntxt}[1]{\text{\textcolor{darkgreen}{#1}}}
\newcommand{\rdtxt}[1]{\text{\textcolor{red}{#1}}}
\newcommand{\code}[1]{\texttt{#1}}
\newcommand{\tldr}{\textbf{TL;DR:}\xspace}
\newcommand{\parahead}[1]{\noindent\textbf{#1.}\enspace}  % run-in heading; use instead of \paragraph to save space
\newcommand*\rot{\rotatebox{90}}

% ---- table cell helpers ----
\newcommand{\cmark}{\ding{51}}
\newcommand{\xmark}{\ding{55}}
\newcommand{\rowhighlight}{\rowcolor{gray!30}}
\newcommand{\cellhighlight}{\cellcolor{yellow!40}}

% ---- rebuttal reviewer tags: bold, distinct colors to signpost which reviewer a point answers ----
\newcommand{\Rone}{\textbf{\textcolor[HTML]{1F77B4}{R1}}}    % blue
\newcommand{\Rtwo}{\textbf{\textcolor[HTML]{D62728}{R2}}}    % red
\newcommand{\Rthree}{\textbf{\textcolor[HTML]{2CA02C}{R3}}}  % green
\newcommand{\Rfour}{\textbf{\textcolor[HTML]{D9730D}{R4}}}   % muted orange

% ---- takeaway callout box ----
\newtcolorbox{takeaway}[1][Takeaway]{
  colback=blue!4!white,      % very light fill
  colframe=blue!50!black,    % muted dark-blue frame
  boxrule=0.6pt,
  arc=2pt,
  left=6pt, right=6pt, top=2pt, bottom=2pt,
  title=#1,
  fonttitle=\bfseries\small,
  coltitle=white,
  colbacktitle=blue!50!black
}

% ==== links & smart references: load LAST (delete if your class already loads hyperref) ====
\usepackage[hidelinks]{hyperref}
\usepackage[capitalize,noabbrev]{cleveref}   % \cref{...}, \Cref{...} — must come after hyperref

% ==== draft mode & \todo ====
% "Explicit" = the main .tex defined \ifdraftmode before \input{packages}.
% (The \ifcsname/\csname dance avoids skipping over an already-defined \ifdraftmode, which
%  crashes the naive \ifdefined...\newif idiom.)
\newif\ifdraftmodeexplicit
\ifcsname ifdraftmode\endcsname \draftmodeexplicittrue \else \draftmodeexplicitfalse \fi
\ifdraftmodeexplicit\else \expandafter\newif\csname ifdraftmode\endcsname \fi

% \todo: visible red box by default and in draft mode; a no-op only when draft mode is
% explicitly OFF (\draftmodefalse in the main .tex), so leftover todos vanish in the final build.
\newcommand{\todo}[1]{\colorbox{red}{TODO: #1}}
\ifdraftmodeexplicit \ifdraftmode\else \renewcommand{\todo}[1]{}\fi \fi

% Draft-mode aids (line numbers, label keys, overfull bars) — on only when draft mode is enabled.
% Worth it for long documents / theses; usually overkill for a short single paper.
% Enable by putting "\newif\ifdraftmode \draftmodetrue" in the main .tex BEFORE \input{packages}.
\ifdraftmode
  \usepackage{lineno}\linenumbers   % line numbers for "see line 214" review comments (prose only; align/equation bodies aren't numbered without extra setup)
  \usepackage{showkeys}             % print \label / \ref keys in the margin
  \overfullrule=5pt                 % black bar in the margin marks every overfull line
\fi
.
% (The \ifcsname/\csname dance avoids skipping over an already-defined \ifdraftmode, which
%  crashes the naive \ifdefined...\newif idiom.)
\newif\ifdraftmodeexplicit
\ifcsname ifdraftmode\endcsname \draftmodeexplicittrue \else \draftmodeexplicitfalse \fi
\ifdraftmodeexplicit\else \expandafter\newif\csname ifdraftmode\endcsname \fi

% \todo: visible red box by default and in draft mode; a no-op only when draft mode is
% explicitly OFF (\draftmodefalse in the main .tex), so leftover todos vanish in the final build.
\newcommand{\todo}[1]{\colorbox{red}{TODO: #1}}
\ifdraftmodeexplicit \ifdraftmode\else \renewcommand{\todo}[1]{}\fi \fi

% Draft-mode aids (line numbers, label keys, overfull bars) — on only when draft mode is enabled.
% Worth it for long documents / theses; usually overkill for a short single paper.
% Enable by putting "\newif\ifdraftmode \draftmodetrue" in the main .tex BEFORE % latex-rules — reusable preamble (see the latex-rules skill).
% Reusable macros only; define paper-specific notation in your main .tex, not here.

% ==== packages ====

% ---- fonts & encoding ----
\usepackage[T1]{fontenc}   % proper 8-bit fonts + real hyphenation
\usepackage{lmodern}       % scalable Latin Modern (fixes bitmap CM warnings)

% ---- math ----
\usepackage{amsmath}
\usepackage{amssymb}
\usepackage{mathtools}     % loads + extends amsmath
\usepackage{dsfont}        % \mathds{1}

% ---- graphics & drawing ----
\usepackage{graphicx}
\usepackage[export]{adjustbox}
\usepackage{tikz}
\usetikzlibrary{shapes}
\usetikzlibrary{positioning}
\usetikzlibrary{trees}
%\usetikzlibrary{graphs,graphdrawing}
%\usegdlibrary{trees}
\usepackage[most]{tcolorbox}
\usepackage{subcaption}    % NOT subfigure / subfig
\usepackage{pdflscape}

% ---- color ----
\usepackage{xcolor}

% ---- tables ----
\usepackage{booktabs}
\usepackage{multirow}
\usepackage{makecell}
\usepackage{colortbl}

% ---- typography, lists, units ----
\usepackage[protrusion=true,expansion=true]{microtype}  % char protrusion + font expansion; drop expansion under XeLaTeX (unsupported)
\usepackage{csquotes}
\usepackage{enumitem}
\usepackage{siunitx}
\sisetup{detect-all, per-mode=symbol}   % detect-all: match surrounding font; per-mode=symbol: m/s not m s^-1
\usepackage{xspace}

% ---- floats, algorithms, symbols ----
\usepackage{float}
\usepackage{afterpage}
\usepackage{footmisc}
\usepackage{listings}
\usepackage{algorithm}
\usepackage{algpseudocode}   % algorithmicx; modern replacement for algorithmic
\usepackage{pifont}
%\usepackage{pgfplots}

% ---- occasionally useful (uncomment as needed) ----
%\usepackage{tabu}
%\usepackage[defaultlines=3,all]{nowidow}
%\usepackage{setspace}
%\usepackage{geometry}
%\usepackage[all]{xy}
%\usepackage{natbib}
%\usepackage{subfigure}
% \usepackage{svg}   % needs --shell-escape + inkscape and is rejected by arXiv; pre-convert SVGs to PDF instead

% ==== math notation ====
% Symbol macros are wrapped in \ensuremath, so they also work in text (\R, \bx, ... auto-enter math).

% ---- quote marks in math ----
\DeclareMathSymbol{\mlq}{\mathord}{operators}{``}
\DeclareMathSymbol{\mrq}{\mathord}{operators}{`'}

% ---- number sets (blackboard bold) ----
% NOTE: \P overrides the built-in pilcrow (¶); see also \S \O \L \H below.
\newcommand{\R}{\ensuremath{\mathbb{R}}}
\newcommand{\N}{\ensuremath{\mathbb{N}}}
\newcommand{\Z}{\ensuremath{\mathbb{Z}}}
\newcommand{\Q}{\ensuremath{\mathbb{Q}}}
\newcommand{\C}{\ensuremath{\mathbb{C}}}
\newcommand{\D}{\ensuremath{\mathbb{D}}}
\newcommand{\E}{\ensuremath{\mathbb{E}}}
\newcommand{\X}{\ensuremath{\mathbb{X}}}
\newcommand{\Y}{\ensuremath{\mathbb{Y}}}
\renewcommand{\P}{\ensuremath{\mathbb{P}}}
\newcommand{\sigB}{\ensuremath{\mathbb{B}}}
\newcommand{\1}{\ensuremath{\mathds{1}}}

% ---- calligraphic sets ----
% NOTE: \S \O \L \H override built-in text commands (§, Ø, Ł, and the o-double-acute accent \H{o}).
\newcommand{\calN}{\ensuremath{\mathcal{N}}}
\newcommand{\calD}{\ensuremath{\mathcal{D}}}
\newcommand{\A}{\ensuremath{\mathcal{A}}}
\newcommand{\B}{\ensuremath{\mathcal{B}}}
\newcommand{\I}{\ensuremath{\mathcal{I}}}
\newcommand{\F}{\ensuremath{\mathcal{F}}}
\newcommand{\G}{\ensuremath{\mathcal{G}}}
\newcommand{\M}{\ensuremath{\mathcal{M}}}
\newcommand{\Tau}{\ensuremath{\mathcal{T}}}
\newcommand{\W}{\ensuremath{\mathcal{W}}}
\newcommand{\sigE}{\ensuremath{\mathcal{E}}}
\renewcommand{\H}{\ensuremath{\mathcal{H}}}
\renewcommand{\L}{\ensuremath{\mathcal{L}}}
\renewcommand{\S}{\ensuremath{\mathcal{S}}}
\renewcommand{\O}{\ensuremath{\mathcal{O}}}

% ---- named operators ----
% NOTE: \Re \Im are redefined to roman "Re"/"Im" (not the default blackletter ℜ/ℑ).
\DeclareMathOperator{\softmax}{softmax}
\DeclareMathOperator*{\argmax}{arg\,max}   % * = subscript sits underneath in display
\DeclareMathOperator{\Var}{Var}
\newcommand{\cov}{\ensuremath{\operatorname{cov}}}
\newcommand{\corr}{\ensuremath{\operatorname{corr}}}
\newcommand{\spanop}{\ensuremath{\operatorname{span}}}
\newcommand{\att}{\ensuremath{\operatorname{attention}}}
\newcommand{\suff}{\ensuremath{\operatorname{suff}}}
\newcommand{\comp}{\ensuremath{\operatorname{comp}}}
\newcommand{\In}{\ensuremath{\operatorname{In}}}
\newcommand{\opcap}{\ensuremath{\operatorname{cap}}}
\newcommand{\entr}{\ensuremath{\operatorname{entries}}}
\newcommand{\ops}{\ensuremath{\operatorname{ops}}}
\newcommand{\id}{\ensuremath{\mathrm{id}}}
\renewcommand{\Re}{\ensuremath{\operatorname{Re}}}
\renewcommand{\Im}{\ensuremath{\operatorname{Im}}}

% ---- symbols & shorthands ----
\newcommand{\eps}{\ensuremath{\varepsilon}}
\renewcommand{\phi}{\ensuremath{\varphi}}     % deliberate: \phi -> varphi
\newcommand{\del}{\ensuremath{\partial}}
\newcommand{\grad}{\ensuremath{\nabla}}
\newcommand{\indep}{\ensuremath{\perp\!\!\!\perp}}
\newcommand{\emptyword}{\ensuremath{\varnothing}}
\newcommand{\tensor}{\ensuremath{\otimes}}
\newcommand{\bigtensor}{\ensuremath{\bigotimes}}
\newcommand{\To}{\ensuremath{\longrightarrow}}
\newcommand{\eqover}[1]{\ensuremath{\stackrel{#1}{=}}}
\newcommand{\darover}[1]{\ensuremath{\xrightarrow{#1}}}
\newcommand{\Hoel}{\text{Höl}}

% ---- paired delimiters ----
\DeclarePairedDelimiter{\norm}{\lVert}{\rVert}   % \norm{x}, \norm*{x} (full size), \norm[\big]{x}
\DeclarePairedDelimiter{\abs}{\lvert}{\rvert}

% ---- bold vectors / matrices ----
\newcommand{\bx}{\ensuremath{\mathbf{x}}}
\newcommand{\bX}{\ensuremath{\mathbf{X}}}
\newcommand{\by}{\ensuremath{\mathbf{y}}}
\newcommand{\bz}{\ensuremath{\mathbf{z}}}
\newcommand{\bB}{\ensuremath{\mathbf{B}}}
\newcommand{\bp}{\ensuremath{\mathbf{p}}}
\newcommand{\bW}{\ensuremath{\mathbf{W}}}
\newcommand{\bt}{\ensuremath{\mathbf{t}}}
\newcommand{\be}{\ensuremath{\mathbf{e}}}

% ---- small constructs ----
\newcommand{\Ball}[2]{\ensuremath{B_#1(#2)}}   % \Ball{\eps}{center}
\newcommand{\twobytwo}[4]{\ensuremath{\begin{pmatrix} #1 & #2 \\ #3 & #4 \end{pmatrix}}}

% ==== text, color & layout helpers ====
\definecolor{darkgreen}{HTML}{009B55}
\newcommand{\gtxt}[1]{\text{\textcolor{gray}{#1}}}
\newcommand{\grntxt}[1]{\text{\textcolor{darkgreen}{#1}}}
\newcommand{\rdtxt}[1]{\text{\textcolor{red}{#1}}}
\newcommand{\code}[1]{\texttt{#1}}
\newcommand{\tldr}{\textbf{TL;DR:}\xspace}
\newcommand{\parahead}[1]{\noindent\textbf{#1.}\enspace}  % run-in heading; use instead of \paragraph to save space
\newcommand*\rot{\rotatebox{90}}

% ---- table cell helpers ----
\newcommand{\cmark}{\ding{51}}
\newcommand{\xmark}{\ding{55}}
\newcommand{\rowhighlight}{\rowcolor{gray!30}}
\newcommand{\cellhighlight}{\cellcolor{yellow!40}}

% ---- rebuttal reviewer tags: bold, distinct colors to signpost which reviewer a point answers ----
\newcommand{\Rone}{\textbf{\textcolor[HTML]{1F77B4}{R1}}}    % blue
\newcommand{\Rtwo}{\textbf{\textcolor[HTML]{D62728}{R2}}}    % red
\newcommand{\Rthree}{\textbf{\textcolor[HTML]{2CA02C}{R3}}}  % green
\newcommand{\Rfour}{\textbf{\textcolor[HTML]{D9730D}{R4}}}   % muted orange

% ---- takeaway callout box ----
\newtcolorbox{takeaway}[1][Takeaway]{
  colback=blue!4!white,      % very light fill
  colframe=blue!50!black,    % muted dark-blue frame
  boxrule=0.6pt,
  arc=2pt,
  left=6pt, right=6pt, top=2pt, bottom=2pt,
  title=#1,
  fonttitle=\bfseries\small,
  coltitle=white,
  colbacktitle=blue!50!black
}

% ==== links & smart references: load LAST (delete if your class already loads hyperref) ====
\usepackage[hidelinks]{hyperref}
\usepackage[capitalize,noabbrev]{cleveref}   % \cref{...}, \Cref{...} — must come after hyperref

% ==== draft mode & \todo ====
% "Explicit" = the main .tex defined \ifdraftmode before \input{packages}.
% (The \ifcsname/\csname dance avoids skipping over an already-defined \ifdraftmode, which
%  crashes the naive \ifdefined...\newif idiom.)
\newif\ifdraftmodeexplicit
\ifcsname ifdraftmode\endcsname \draftmodeexplicittrue \else \draftmodeexplicitfalse \fi
\ifdraftmodeexplicit\else \expandafter\newif\csname ifdraftmode\endcsname \fi

% \todo: visible red box by default and in draft mode; a no-op only when draft mode is
% explicitly OFF (\draftmodefalse in the main .tex), so leftover todos vanish in the final build.
\newcommand{\todo}[1]{\colorbox{red}{TODO: #1}}
\ifdraftmodeexplicit \ifdraftmode\else \renewcommand{\todo}[1]{}\fi \fi

% Draft-mode aids (line numbers, label keys, overfull bars) — on only when draft mode is enabled.
% Worth it for long documents / theses; usually overkill for a short single paper.
% Enable by putting "\newif\ifdraftmode \draftmodetrue" in the main .tex BEFORE \input{packages}.
\ifdraftmode
  \usepackage{lineno}\linenumbers   % line numbers for "see line 214" review comments (prose only; align/equation bodies aren't numbered without extra setup)
  \usepackage{showkeys}             % print \label / \ref keys in the margin
  \overfullrule=5pt                 % black bar in the margin marks every overfull line
\fi
.
\ifdraftmode
  \usepackage{lineno}\linenumbers   % line numbers for "see line 214" review comments (prose only; align/equation bodies aren't numbered without extra setup)
  \usepackage{showkeys}             % print \label / \ref keys in the margin
  \overfullrule=5pt                 % black bar in the margin marks every overfull line
\fi
.
% (The \ifcsname/\csname dance avoids skipping over an already-defined \ifdraftmode, which
%  crashes the naive \ifdefined...\newif idiom.)
\newif\ifdraftmodeexplicit
\ifcsname ifdraftmode\endcsname \draftmodeexplicittrue \else \draftmodeexplicitfalse \fi
\ifdraftmodeexplicit\else \expandafter\newif\csname ifdraftmode\endcsname \fi

% \todo: visible red box by default and in draft mode; a no-op only when draft mode is
% explicitly OFF (\draftmodefalse in the main .tex), so leftover todos vanish in the final build.
\newcommand{\todo}[1]{\colorbox{red}{TODO: #1}}
\ifdraftmodeexplicit \ifdraftmode\else \renewcommand{\todo}[1]{}\fi \fi

% Draft-mode aids (line numbers, label keys, overfull bars) — on only when draft mode is enabled.
% Worth it for long documents / theses; usually overkill for a short single paper.
% Enable by putting "\newif\ifdraftmode \draftmodetrue" in the main .tex BEFORE % latex-rules — reusable preamble (see the latex-rules skill).
% Reusable macros only; define paper-specific notation in your main .tex, not here.

% ==== packages ====

% ---- fonts & encoding ----
\usepackage[T1]{fontenc}   % proper 8-bit fonts + real hyphenation
\usepackage{lmodern}       % scalable Latin Modern (fixes bitmap CM warnings)

% ---- math ----
\usepackage{amsmath}
\usepackage{amssymb}
\usepackage{mathtools}     % loads + extends amsmath
\usepackage{dsfont}        % \mathds{1}

% ---- graphics & drawing ----
\usepackage{graphicx}
\usepackage[export]{adjustbox}
\usepackage{tikz}
\usetikzlibrary{shapes}
\usetikzlibrary{positioning}
\usetikzlibrary{trees}
%\usetikzlibrary{graphs,graphdrawing}
%\usegdlibrary{trees}
\usepackage[most]{tcolorbox}
\usepackage{subcaption}    % NOT subfigure / subfig
\usepackage{pdflscape}

% ---- color ----
\usepackage{xcolor}

% ---- tables ----
\usepackage{booktabs}
\usepackage{multirow}
\usepackage{makecell}
\usepackage{colortbl}

% ---- typography, lists, units ----
\usepackage[protrusion=true,expansion=true]{microtype}  % char protrusion + font expansion; drop expansion under XeLaTeX (unsupported)
\usepackage{csquotes}
\usepackage{enumitem}
\usepackage{siunitx}
\sisetup{detect-all, per-mode=symbol}   % detect-all: match surrounding font; per-mode=symbol: m/s not m s^-1
\usepackage{xspace}

% ---- floats, algorithms, symbols ----
\usepackage{float}
\usepackage{afterpage}
\usepackage{footmisc}
\usepackage{listings}
\usepackage{algorithm}
\usepackage{algpseudocode}   % algorithmicx; modern replacement for algorithmic
\usepackage{pifont}
%\usepackage{pgfplots}

% ---- occasionally useful (uncomment as needed) ----
%\usepackage{tabu}
%\usepackage[defaultlines=3,all]{nowidow}
%\usepackage{setspace}
%\usepackage{geometry}
%\usepackage[all]{xy}
%\usepackage{natbib}
%\usepackage{subfigure}
% \usepackage{svg}   % needs --shell-escape + inkscape and is rejected by arXiv; pre-convert SVGs to PDF instead

% ==== math notation ====
% Symbol macros are wrapped in \ensuremath, so they also work in text (\R, \bx, ... auto-enter math).

% ---- quote marks in math ----
\DeclareMathSymbol{\mlq}{\mathord}{operators}{``}
\DeclareMathSymbol{\mrq}{\mathord}{operators}{`'}

% ---- number sets (blackboard bold) ----
% NOTE: \P overrides the built-in pilcrow (¶); see also \S \O \L \H below.
\newcommand{\R}{\ensuremath{\mathbb{R}}}
\newcommand{\N}{\ensuremath{\mathbb{N}}}
\newcommand{\Z}{\ensuremath{\mathbb{Z}}}
\newcommand{\Q}{\ensuremath{\mathbb{Q}}}
\newcommand{\C}{\ensuremath{\mathbb{C}}}
\newcommand{\D}{\ensuremath{\mathbb{D}}}
\newcommand{\E}{\ensuremath{\mathbb{E}}}
\newcommand{\X}{\ensuremath{\mathbb{X}}}
\newcommand{\Y}{\ensuremath{\mathbb{Y}}}
\renewcommand{\P}{\ensuremath{\mathbb{P}}}
\newcommand{\sigB}{\ensuremath{\mathbb{B}}}
\newcommand{\1}{\ensuremath{\mathds{1}}}

% ---- calligraphic sets ----
% NOTE: \S \O \L \H override built-in text commands (§, Ø, Ł, and the o-double-acute accent \H{o}).
\newcommand{\calN}{\ensuremath{\mathcal{N}}}
\newcommand{\calD}{\ensuremath{\mathcal{D}}}
\newcommand{\A}{\ensuremath{\mathcal{A}}}
\newcommand{\B}{\ensuremath{\mathcal{B}}}
\newcommand{\I}{\ensuremath{\mathcal{I}}}
\newcommand{\F}{\ensuremath{\mathcal{F}}}
\newcommand{\G}{\ensuremath{\mathcal{G}}}
\newcommand{\M}{\ensuremath{\mathcal{M}}}
\newcommand{\Tau}{\ensuremath{\mathcal{T}}}
\newcommand{\W}{\ensuremath{\mathcal{W}}}
\newcommand{\sigE}{\ensuremath{\mathcal{E}}}
\renewcommand{\H}{\ensuremath{\mathcal{H}}}
\renewcommand{\L}{\ensuremath{\mathcal{L}}}
\renewcommand{\S}{\ensuremath{\mathcal{S}}}
\renewcommand{\O}{\ensuremath{\mathcal{O}}}

% ---- named operators ----
% NOTE: \Re \Im are redefined to roman "Re"/"Im" (not the default blackletter ℜ/ℑ).
\DeclareMathOperator{\softmax}{softmax}
\DeclareMathOperator*{\argmax}{arg\,max}   % * = subscript sits underneath in display
\DeclareMathOperator{\Var}{Var}
\newcommand{\cov}{\ensuremath{\operatorname{cov}}}
\newcommand{\corr}{\ensuremath{\operatorname{corr}}}
\newcommand{\spanop}{\ensuremath{\operatorname{span}}}
\newcommand{\att}{\ensuremath{\operatorname{attention}}}
\newcommand{\suff}{\ensuremath{\operatorname{suff}}}
\newcommand{\comp}{\ensuremath{\operatorname{comp}}}
\newcommand{\In}{\ensuremath{\operatorname{In}}}
\newcommand{\opcap}{\ensuremath{\operatorname{cap}}}
\newcommand{\entr}{\ensuremath{\operatorname{entries}}}
\newcommand{\ops}{\ensuremath{\operatorname{ops}}}
\newcommand{\id}{\ensuremath{\mathrm{id}}}
\renewcommand{\Re}{\ensuremath{\operatorname{Re}}}
\renewcommand{\Im}{\ensuremath{\operatorname{Im}}}

% ---- symbols & shorthands ----
\newcommand{\eps}{\ensuremath{\varepsilon}}
\renewcommand{\phi}{\ensuremath{\varphi}}     % deliberate: \phi -> varphi
\newcommand{\del}{\ensuremath{\partial}}
\newcommand{\grad}{\ensuremath{\nabla}}
\newcommand{\indep}{\ensuremath{\perp\!\!\!\perp}}
\newcommand{\emptyword}{\ensuremath{\varnothing}}
\newcommand{\tensor}{\ensuremath{\otimes}}
\newcommand{\bigtensor}{\ensuremath{\bigotimes}}
\newcommand{\To}{\ensuremath{\longrightarrow}}
\newcommand{\eqover}[1]{\ensuremath{\stackrel{#1}{=}}}
\newcommand{\darover}[1]{\ensuremath{\xrightarrow{#1}}}
\newcommand{\Hoel}{\text{Höl}}

% ---- paired delimiters ----
\DeclarePairedDelimiter{\norm}{\lVert}{\rVert}   % \norm{x}, \norm*{x} (full size), \norm[\big]{x}
\DeclarePairedDelimiter{\abs}{\lvert}{\rvert}

% ---- bold vectors / matrices ----
\newcommand{\bx}{\ensuremath{\mathbf{x}}}
\newcommand{\bX}{\ensuremath{\mathbf{X}}}
\newcommand{\by}{\ensuremath{\mathbf{y}}}
\newcommand{\bz}{\ensuremath{\mathbf{z}}}
\newcommand{\bB}{\ensuremath{\mathbf{B}}}
\newcommand{\bp}{\ensuremath{\mathbf{p}}}
\newcommand{\bW}{\ensuremath{\mathbf{W}}}
\newcommand{\bt}{\ensuremath{\mathbf{t}}}
\newcommand{\be}{\ensuremath{\mathbf{e}}}

% ---- small constructs ----
\newcommand{\Ball}[2]{\ensuremath{B_#1(#2)}}   % \Ball{\eps}{center}
\newcommand{\twobytwo}[4]{\ensuremath{\begin{pmatrix} #1 & #2 \\ #3 & #4 \end{pmatrix}}}

% ==== text, color & layout helpers ====
\definecolor{darkgreen}{HTML}{009B55}
\newcommand{\gtxt}[1]{\text{\textcolor{gray}{#1}}}
\newcommand{\grntxt}[1]{\text{\textcolor{darkgreen}{#1}}}
\newcommand{\rdtxt}[1]{\text{\textcolor{red}{#1}}}
\newcommand{\code}[1]{\texttt{#1}}
\newcommand{\tldr}{\textbf{TL;DR:}\xspace}
\newcommand{\parahead}[1]{\noindent\textbf{#1.}\enspace}  % run-in heading; use instead of \paragraph to save space
\newcommand*\rot{\rotatebox{90}}

% ---- table cell helpers ----
\newcommand{\cmark}{\ding{51}}
\newcommand{\xmark}{\ding{55}}
\newcommand{\rowhighlight}{\rowcolor{gray!30}}
\newcommand{\cellhighlight}{\cellcolor{yellow!40}}

% ---- rebuttal reviewer tags: bold, distinct colors to signpost which reviewer a point answers ----
\newcommand{\Rone}{\textbf{\textcolor[HTML]{1F77B4}{R1}}}    % blue
\newcommand{\Rtwo}{\textbf{\textcolor[HTML]{D62728}{R2}}}    % red
\newcommand{\Rthree}{\textbf{\textcolor[HTML]{2CA02C}{R3}}}  % green
\newcommand{\Rfour}{\textbf{\textcolor[HTML]{D9730D}{R4}}}   % muted orange

% ---- takeaway callout box ----
\newtcolorbox{takeaway}[1][Takeaway]{
  colback=blue!4!white,      % very light fill
  colframe=blue!50!black,    % muted dark-blue frame
  boxrule=0.6pt,
  arc=2pt,
  left=6pt, right=6pt, top=2pt, bottom=2pt,
  title=#1,
  fonttitle=\bfseries\small,
  coltitle=white,
  colbacktitle=blue!50!black
}

% ==== links & smart references: load LAST (delete if your class already loads hyperref) ====
\usepackage[hidelinks]{hyperref}
\usepackage[capitalize,noabbrev]{cleveref}   % \cref{...}, \Cref{...} — must come after hyperref

% ==== draft mode & \todo ====
% "Explicit" = the main .tex defined \ifdraftmode before % latex-rules — reusable preamble (see the latex-rules skill).
% Reusable macros only; define paper-specific notation in your main .tex, not here.

% ==== packages ====

% ---- fonts & encoding ----
\usepackage[T1]{fontenc}   % proper 8-bit fonts + real hyphenation
\usepackage{lmodern}       % scalable Latin Modern (fixes bitmap CM warnings)

% ---- math ----
\usepackage{amsmath}
\usepackage{amssymb}
\usepackage{mathtools}     % loads + extends amsmath
\usepackage{dsfont}        % \mathds{1}

% ---- graphics & drawing ----
\usepackage{graphicx}
\usepackage[export]{adjustbox}
\usepackage{tikz}
\usetikzlibrary{shapes}
\usetikzlibrary{positioning}
\usetikzlibrary{trees}
%\usetikzlibrary{graphs,graphdrawing}
%\usegdlibrary{trees}
\usepackage[most]{tcolorbox}
\usepackage{subcaption}    % NOT subfigure / subfig
\usepackage{pdflscape}

% ---- color ----
\usepackage{xcolor}

% ---- tables ----
\usepackage{booktabs}
\usepackage{multirow}
\usepackage{makecell}
\usepackage{colortbl}

% ---- typography, lists, units ----
\usepackage[protrusion=true,expansion=true]{microtype}  % char protrusion + font expansion; drop expansion under XeLaTeX (unsupported)
\usepackage{csquotes}
\usepackage{enumitem}
\usepackage{siunitx}
\sisetup{detect-all, per-mode=symbol}   % detect-all: match surrounding font; per-mode=symbol: m/s not m s^-1
\usepackage{xspace}

% ---- floats, algorithms, symbols ----
\usepackage{float}
\usepackage{afterpage}
\usepackage{footmisc}
\usepackage{listings}
\usepackage{algorithm}
\usepackage{algpseudocode}   % algorithmicx; modern replacement for algorithmic
\usepackage{pifont}
%\usepackage{pgfplots}

% ---- occasionally useful (uncomment as needed) ----
%\usepackage{tabu}
%\usepackage[defaultlines=3,all]{nowidow}
%\usepackage{setspace}
%\usepackage{geometry}
%\usepackage[all]{xy}
%\usepackage{natbib}
%\usepackage{subfigure}
% \usepackage{svg}   % needs --shell-escape + inkscape and is rejected by arXiv; pre-convert SVGs to PDF instead

% ==== math notation ====
% Symbol macros are wrapped in \ensuremath, so they also work in text (\R, \bx, ... auto-enter math).

% ---- quote marks in math ----
\DeclareMathSymbol{\mlq}{\mathord}{operators}{``}
\DeclareMathSymbol{\mrq}{\mathord}{operators}{`'}

% ---- number sets (blackboard bold) ----
% NOTE: \P overrides the built-in pilcrow (¶); see also \S \O \L \H below.
\newcommand{\R}{\ensuremath{\mathbb{R}}}
\newcommand{\N}{\ensuremath{\mathbb{N}}}
\newcommand{\Z}{\ensuremath{\mathbb{Z}}}
\newcommand{\Q}{\ensuremath{\mathbb{Q}}}
\newcommand{\C}{\ensuremath{\mathbb{C}}}
\newcommand{\D}{\ensuremath{\mathbb{D}}}
\newcommand{\E}{\ensuremath{\mathbb{E}}}
\newcommand{\X}{\ensuremath{\mathbb{X}}}
\newcommand{\Y}{\ensuremath{\mathbb{Y}}}
\renewcommand{\P}{\ensuremath{\mathbb{P}}}
\newcommand{\sigB}{\ensuremath{\mathbb{B}}}
\newcommand{\1}{\ensuremath{\mathds{1}}}

% ---- calligraphic sets ----
% NOTE: \S \O \L \H override built-in text commands (§, Ø, Ł, and the o-double-acute accent \H{o}).
\newcommand{\calN}{\ensuremath{\mathcal{N}}}
\newcommand{\calD}{\ensuremath{\mathcal{D}}}
\newcommand{\A}{\ensuremath{\mathcal{A}}}
\newcommand{\B}{\ensuremath{\mathcal{B}}}
\newcommand{\I}{\ensuremath{\mathcal{I}}}
\newcommand{\F}{\ensuremath{\mathcal{F}}}
\newcommand{\G}{\ensuremath{\mathcal{G}}}
\newcommand{\M}{\ensuremath{\mathcal{M}}}
\newcommand{\Tau}{\ensuremath{\mathcal{T}}}
\newcommand{\W}{\ensuremath{\mathcal{W}}}
\newcommand{\sigE}{\ensuremath{\mathcal{E}}}
\renewcommand{\H}{\ensuremath{\mathcal{H}}}
\renewcommand{\L}{\ensuremath{\mathcal{L}}}
\renewcommand{\S}{\ensuremath{\mathcal{S}}}
\renewcommand{\O}{\ensuremath{\mathcal{O}}}

% ---- named operators ----
% NOTE: \Re \Im are redefined to roman "Re"/"Im" (not the default blackletter ℜ/ℑ).
\DeclareMathOperator{\softmax}{softmax}
\DeclareMathOperator*{\argmax}{arg\,max}   % * = subscript sits underneath in display
\DeclareMathOperator{\Var}{Var}
\newcommand{\cov}{\ensuremath{\operatorname{cov}}}
\newcommand{\corr}{\ensuremath{\operatorname{corr}}}
\newcommand{\spanop}{\ensuremath{\operatorname{span}}}
\newcommand{\att}{\ensuremath{\operatorname{attention}}}
\newcommand{\suff}{\ensuremath{\operatorname{suff}}}
\newcommand{\comp}{\ensuremath{\operatorname{comp}}}
\newcommand{\In}{\ensuremath{\operatorname{In}}}
\newcommand{\opcap}{\ensuremath{\operatorname{cap}}}
\newcommand{\entr}{\ensuremath{\operatorname{entries}}}
\newcommand{\ops}{\ensuremath{\operatorname{ops}}}
\newcommand{\id}{\ensuremath{\mathrm{id}}}
\renewcommand{\Re}{\ensuremath{\operatorname{Re}}}
\renewcommand{\Im}{\ensuremath{\operatorname{Im}}}

% ---- symbols & shorthands ----
\newcommand{\eps}{\ensuremath{\varepsilon}}
\renewcommand{\phi}{\ensuremath{\varphi}}     % deliberate: \phi -> varphi
\newcommand{\del}{\ensuremath{\partial}}
\newcommand{\grad}{\ensuremath{\nabla}}
\newcommand{\indep}{\ensuremath{\perp\!\!\!\perp}}
\newcommand{\emptyword}{\ensuremath{\varnothing}}
\newcommand{\tensor}{\ensuremath{\otimes}}
\newcommand{\bigtensor}{\ensuremath{\bigotimes}}
\newcommand{\To}{\ensuremath{\longrightarrow}}
\newcommand{\eqover}[1]{\ensuremath{\stackrel{#1}{=}}}
\newcommand{\darover}[1]{\ensuremath{\xrightarrow{#1}}}
\newcommand{\Hoel}{\text{Höl}}

% ---- paired delimiters ----
\DeclarePairedDelimiter{\norm}{\lVert}{\rVert}   % \norm{x}, \norm*{x} (full size), \norm[\big]{x}
\DeclarePairedDelimiter{\abs}{\lvert}{\rvert}

% ---- bold vectors / matrices ----
\newcommand{\bx}{\ensuremath{\mathbf{x}}}
\newcommand{\bX}{\ensuremath{\mathbf{X}}}
\newcommand{\by}{\ensuremath{\mathbf{y}}}
\newcommand{\bz}{\ensuremath{\mathbf{z}}}
\newcommand{\bB}{\ensuremath{\mathbf{B}}}
\newcommand{\bp}{\ensuremath{\mathbf{p}}}
\newcommand{\bW}{\ensuremath{\mathbf{W}}}
\newcommand{\bt}{\ensuremath{\mathbf{t}}}
\newcommand{\be}{\ensuremath{\mathbf{e}}}

% ---- small constructs ----
\newcommand{\Ball}[2]{\ensuremath{B_#1(#2)}}   % \Ball{\eps}{center}
\newcommand{\twobytwo}[4]{\ensuremath{\begin{pmatrix} #1 & #2 \\ #3 & #4 \end{pmatrix}}}

% ==== text, color & layout helpers ====
\definecolor{darkgreen}{HTML}{009B55}
\newcommand{\gtxt}[1]{\text{\textcolor{gray}{#1}}}
\newcommand{\grntxt}[1]{\text{\textcolor{darkgreen}{#1}}}
\newcommand{\rdtxt}[1]{\text{\textcolor{red}{#1}}}
\newcommand{\code}[1]{\texttt{#1}}
\newcommand{\tldr}{\textbf{TL;DR:}\xspace}
\newcommand{\parahead}[1]{\noindent\textbf{#1.}\enspace}  % run-in heading; use instead of \paragraph to save space
\newcommand*\rot{\rotatebox{90}}

% ---- table cell helpers ----
\newcommand{\cmark}{\ding{51}}
\newcommand{\xmark}{\ding{55}}
\newcommand{\rowhighlight}{\rowcolor{gray!30}}
\newcommand{\cellhighlight}{\cellcolor{yellow!40}}

% ---- rebuttal reviewer tags: bold, distinct colors to signpost which reviewer a point answers ----
\newcommand{\Rone}{\textbf{\textcolor[HTML]{1F77B4}{R1}}}    % blue
\newcommand{\Rtwo}{\textbf{\textcolor[HTML]{D62728}{R2}}}    % red
\newcommand{\Rthree}{\textbf{\textcolor[HTML]{2CA02C}{R3}}}  % green
\newcommand{\Rfour}{\textbf{\textcolor[HTML]{D9730D}{R4}}}   % muted orange

% ---- takeaway callout box ----
\newtcolorbox{takeaway}[1][Takeaway]{
  colback=blue!4!white,      % very light fill
  colframe=blue!50!black,    % muted dark-blue frame
  boxrule=0.6pt,
  arc=2pt,
  left=6pt, right=6pt, top=2pt, bottom=2pt,
  title=#1,
  fonttitle=\bfseries\small,
  coltitle=white,
  colbacktitle=blue!50!black
}

% ==== links & smart references: load LAST (delete if your class already loads hyperref) ====
\usepackage[hidelinks]{hyperref}
\usepackage[capitalize,noabbrev]{cleveref}   % \cref{...}, \Cref{...} — must come after hyperref

% ==== draft mode & \todo ====
% "Explicit" = the main .tex defined \ifdraftmode before \input{packages}.
% (The \ifcsname/\csname dance avoids skipping over an already-defined \ifdraftmode, which
%  crashes the naive \ifdefined...\newif idiom.)
\newif\ifdraftmodeexplicit
\ifcsname ifdraftmode\endcsname \draftmodeexplicittrue \else \draftmodeexplicitfalse \fi
\ifdraftmodeexplicit\else \expandafter\newif\csname ifdraftmode\endcsname \fi

% \todo: visible red box by default and in draft mode; a no-op only when draft mode is
% explicitly OFF (\draftmodefalse in the main .tex), so leftover todos vanish in the final build.
\newcommand{\todo}[1]{\colorbox{red}{TODO: #1}}
\ifdraftmodeexplicit \ifdraftmode\else \renewcommand{\todo}[1]{}\fi \fi

% Draft-mode aids (line numbers, label keys, overfull bars) — on only when draft mode is enabled.
% Worth it for long documents / theses; usually overkill for a short single paper.
% Enable by putting "\newif\ifdraftmode \draftmodetrue" in the main .tex BEFORE \input{packages}.
\ifdraftmode
  \usepackage{lineno}\linenumbers   % line numbers for "see line 214" review comments (prose only; align/equation bodies aren't numbered without extra setup)
  \usepackage{showkeys}             % print \label / \ref keys in the margin
  \overfullrule=5pt                 % black bar in the margin marks every overfull line
\fi
.
% (The \ifcsname/\csname dance avoids skipping over an already-defined \ifdraftmode, which
%  crashes the naive \ifdefined...\newif idiom.)
\newif\ifdraftmodeexplicit
\ifcsname ifdraftmode\endcsname \draftmodeexplicittrue \else \draftmodeexplicitfalse \fi
\ifdraftmodeexplicit\else \expandafter\newif\csname ifdraftmode\endcsname \fi

% \todo: visible red box by default and in draft mode; a no-op only when draft mode is
% explicitly OFF (\draftmodefalse in the main .tex), so leftover todos vanish in the final build.
\newcommand{\todo}[1]{\colorbox{red}{TODO: #1}}
\ifdraftmodeexplicit \ifdraftmode\else \renewcommand{\todo}[1]{}\fi \fi

% Draft-mode aids (line numbers, label keys, overfull bars) — on only when draft mode is enabled.
% Worth it for long documents / theses; usually overkill for a short single paper.
% Enable by putting "\newif\ifdraftmode \draftmodetrue" in the main .tex BEFORE % latex-rules — reusable preamble (see the latex-rules skill).
% Reusable macros only; define paper-specific notation in your main .tex, not here.

% ==== packages ====

% ---- fonts & encoding ----
\usepackage[T1]{fontenc}   % proper 8-bit fonts + real hyphenation
\usepackage{lmodern}       % scalable Latin Modern (fixes bitmap CM warnings)

% ---- math ----
\usepackage{amsmath}
\usepackage{amssymb}
\usepackage{mathtools}     % loads + extends amsmath
\usepackage{dsfont}        % \mathds{1}

% ---- graphics & drawing ----
\usepackage{graphicx}
\usepackage[export]{adjustbox}
\usepackage{tikz}
\usetikzlibrary{shapes}
\usetikzlibrary{positioning}
\usetikzlibrary{trees}
%\usetikzlibrary{graphs,graphdrawing}
%\usegdlibrary{trees}
\usepackage[most]{tcolorbox}
\usepackage{subcaption}    % NOT subfigure / subfig
\usepackage{pdflscape}

% ---- color ----
\usepackage{xcolor}

% ---- tables ----
\usepackage{booktabs}
\usepackage{multirow}
\usepackage{makecell}
\usepackage{colortbl}

% ---- typography, lists, units ----
\usepackage[protrusion=true,expansion=true]{microtype}  % char protrusion + font expansion; drop expansion under XeLaTeX (unsupported)
\usepackage{csquotes}
\usepackage{enumitem}
\usepackage{siunitx}
\sisetup{detect-all, per-mode=symbol}   % detect-all: match surrounding font; per-mode=symbol: m/s not m s^-1
\usepackage{xspace}

% ---- floats, algorithms, symbols ----
\usepackage{float}
\usepackage{afterpage}
\usepackage{footmisc}
\usepackage{listings}
\usepackage{algorithm}
\usepackage{algpseudocode}   % algorithmicx; modern replacement for algorithmic
\usepackage{pifont}
%\usepackage{pgfplots}

% ---- occasionally useful (uncomment as needed) ----
%\usepackage{tabu}
%\usepackage[defaultlines=3,all]{nowidow}
%\usepackage{setspace}
%\usepackage{geometry}
%\usepackage[all]{xy}
%\usepackage{natbib}
%\usepackage{subfigure}
% \usepackage{svg}   % needs --shell-escape + inkscape and is rejected by arXiv; pre-convert SVGs to PDF instead

% ==== math notation ====
% Symbol macros are wrapped in \ensuremath, so they also work in text (\R, \bx, ... auto-enter math).

% ---- quote marks in math ----
\DeclareMathSymbol{\mlq}{\mathord}{operators}{``}
\DeclareMathSymbol{\mrq}{\mathord}{operators}{`'}

% ---- number sets (blackboard bold) ----
% NOTE: \P overrides the built-in pilcrow (¶); see also \S \O \L \H below.
\newcommand{\R}{\ensuremath{\mathbb{R}}}
\newcommand{\N}{\ensuremath{\mathbb{N}}}
\newcommand{\Z}{\ensuremath{\mathbb{Z}}}
\newcommand{\Q}{\ensuremath{\mathbb{Q}}}
\newcommand{\C}{\ensuremath{\mathbb{C}}}
\newcommand{\D}{\ensuremath{\mathbb{D}}}
\newcommand{\E}{\ensuremath{\mathbb{E}}}
\newcommand{\X}{\ensuremath{\mathbb{X}}}
\newcommand{\Y}{\ensuremath{\mathbb{Y}}}
\renewcommand{\P}{\ensuremath{\mathbb{P}}}
\newcommand{\sigB}{\ensuremath{\mathbb{B}}}
\newcommand{\1}{\ensuremath{\mathds{1}}}

% ---- calligraphic sets ----
% NOTE: \S \O \L \H override built-in text commands (§, Ø, Ł, and the o-double-acute accent \H{o}).
\newcommand{\calN}{\ensuremath{\mathcal{N}}}
\newcommand{\calD}{\ensuremath{\mathcal{D}}}
\newcommand{\A}{\ensuremath{\mathcal{A}}}
\newcommand{\B}{\ensuremath{\mathcal{B}}}
\newcommand{\I}{\ensuremath{\mathcal{I}}}
\newcommand{\F}{\ensuremath{\mathcal{F}}}
\newcommand{\G}{\ensuremath{\mathcal{G}}}
\newcommand{\M}{\ensuremath{\mathcal{M}}}
\newcommand{\Tau}{\ensuremath{\mathcal{T}}}
\newcommand{\W}{\ensuremath{\mathcal{W}}}
\newcommand{\sigE}{\ensuremath{\mathcal{E}}}
\renewcommand{\H}{\ensuremath{\mathcal{H}}}
\renewcommand{\L}{\ensuremath{\mathcal{L}}}
\renewcommand{\S}{\ensuremath{\mathcal{S}}}
\renewcommand{\O}{\ensuremath{\mathcal{O}}}

% ---- named operators ----
% NOTE: \Re \Im are redefined to roman "Re"/"Im" (not the default blackletter ℜ/ℑ).
\DeclareMathOperator{\softmax}{softmax}
\DeclareMathOperator*{\argmax}{arg\,max}   % * = subscript sits underneath in display
\DeclareMathOperator{\Var}{Var}
\newcommand{\cov}{\ensuremath{\operatorname{cov}}}
\newcommand{\corr}{\ensuremath{\operatorname{corr}}}
\newcommand{\spanop}{\ensuremath{\operatorname{span}}}
\newcommand{\att}{\ensuremath{\operatorname{attention}}}
\newcommand{\suff}{\ensuremath{\operatorname{suff}}}
\newcommand{\comp}{\ensuremath{\operatorname{comp}}}
\newcommand{\In}{\ensuremath{\operatorname{In}}}
\newcommand{\opcap}{\ensuremath{\operatorname{cap}}}
\newcommand{\entr}{\ensuremath{\operatorname{entries}}}
\newcommand{\ops}{\ensuremath{\operatorname{ops}}}
\newcommand{\id}{\ensuremath{\mathrm{id}}}
\renewcommand{\Re}{\ensuremath{\operatorname{Re}}}
\renewcommand{\Im}{\ensuremath{\operatorname{Im}}}

% ---- symbols & shorthands ----
\newcommand{\eps}{\ensuremath{\varepsilon}}
\renewcommand{\phi}{\ensuremath{\varphi}}     % deliberate: \phi -> varphi
\newcommand{\del}{\ensuremath{\partial}}
\newcommand{\grad}{\ensuremath{\nabla}}
\newcommand{\indep}{\ensuremath{\perp\!\!\!\perp}}
\newcommand{\emptyword}{\ensuremath{\varnothing}}
\newcommand{\tensor}{\ensuremath{\otimes}}
\newcommand{\bigtensor}{\ensuremath{\bigotimes}}
\newcommand{\To}{\ensuremath{\longrightarrow}}
\newcommand{\eqover}[1]{\ensuremath{\stackrel{#1}{=}}}
\newcommand{\darover}[1]{\ensuremath{\xrightarrow{#1}}}
\newcommand{\Hoel}{\text{Höl}}

% ---- paired delimiters ----
\DeclarePairedDelimiter{\norm}{\lVert}{\rVert}   % \norm{x}, \norm*{x} (full size), \norm[\big]{x}
\DeclarePairedDelimiter{\abs}{\lvert}{\rvert}

% ---- bold vectors / matrices ----
\newcommand{\bx}{\ensuremath{\mathbf{x}}}
\newcommand{\bX}{\ensuremath{\mathbf{X}}}
\newcommand{\by}{\ensuremath{\mathbf{y}}}
\newcommand{\bz}{\ensuremath{\mathbf{z}}}
\newcommand{\bB}{\ensuremath{\mathbf{B}}}
\newcommand{\bp}{\ensuremath{\mathbf{p}}}
\newcommand{\bW}{\ensuremath{\mathbf{W}}}
\newcommand{\bt}{\ensuremath{\mathbf{t}}}
\newcommand{\be}{\ensuremath{\mathbf{e}}}

% ---- small constructs ----
\newcommand{\Ball}[2]{\ensuremath{B_#1(#2)}}   % \Ball{\eps}{center}
\newcommand{\twobytwo}[4]{\ensuremath{\begin{pmatrix} #1 & #2 \\ #3 & #4 \end{pmatrix}}}

% ==== text, color & layout helpers ====
\definecolor{darkgreen}{HTML}{009B55}
\newcommand{\gtxt}[1]{\text{\textcolor{gray}{#1}}}
\newcommand{\grntxt}[1]{\text{\textcolor{darkgreen}{#1}}}
\newcommand{\rdtxt}[1]{\text{\textcolor{red}{#1}}}
\newcommand{\code}[1]{\texttt{#1}}
\newcommand{\tldr}{\textbf{TL;DR:}\xspace}
\newcommand{\parahead}[1]{\noindent\textbf{#1.}\enspace}  % run-in heading; use instead of \paragraph to save space
\newcommand*\rot{\rotatebox{90}}

% ---- table cell helpers ----
\newcommand{\cmark}{\ding{51}}
\newcommand{\xmark}{\ding{55}}
\newcommand{\rowhighlight}{\rowcolor{gray!30}}
\newcommand{\cellhighlight}{\cellcolor{yellow!40}}

% ---- rebuttal reviewer tags: bold, distinct colors to signpost which reviewer a point answers ----
\newcommand{\Rone}{\textbf{\textcolor[HTML]{1F77B4}{R1}}}    % blue
\newcommand{\Rtwo}{\textbf{\textcolor[HTML]{D62728}{R2}}}    % red
\newcommand{\Rthree}{\textbf{\textcolor[HTML]{2CA02C}{R3}}}  % green
\newcommand{\Rfour}{\textbf{\textcolor[HTML]{D9730D}{R4}}}   % muted orange

% ---- takeaway callout box ----
\newtcolorbox{takeaway}[1][Takeaway]{
  colback=blue!4!white,      % very light fill
  colframe=blue!50!black,    % muted dark-blue frame
  boxrule=0.6pt,
  arc=2pt,
  left=6pt, right=6pt, top=2pt, bottom=2pt,
  title=#1,
  fonttitle=\bfseries\small,
  coltitle=white,
  colbacktitle=blue!50!black
}

% ==== links & smart references: load LAST (delete if your class already loads hyperref) ====
\usepackage[hidelinks]{hyperref}
\usepackage[capitalize,noabbrev]{cleveref}   % \cref{...}, \Cref{...} — must come after hyperref

% ==== draft mode & \todo ====
% "Explicit" = the main .tex defined \ifdraftmode before \input{packages}.
% (The \ifcsname/\csname dance avoids skipping over an already-defined \ifdraftmode, which
%  crashes the naive \ifdefined...\newif idiom.)
\newif\ifdraftmodeexplicit
\ifcsname ifdraftmode\endcsname \draftmodeexplicittrue \else \draftmodeexplicitfalse \fi
\ifdraftmodeexplicit\else \expandafter\newif\csname ifdraftmode\endcsname \fi

% \todo: visible red box by default and in draft mode; a no-op only when draft mode is
% explicitly OFF (\draftmodefalse in the main .tex), so leftover todos vanish in the final build.
\newcommand{\todo}[1]{\colorbox{red}{TODO: #1}}
\ifdraftmodeexplicit \ifdraftmode\else \renewcommand{\todo}[1]{}\fi \fi

% Draft-mode aids (line numbers, label keys, overfull bars) — on only when draft mode is enabled.
% Worth it for long documents / theses; usually overkill for a short single paper.
% Enable by putting "\newif\ifdraftmode \draftmodetrue" in the main .tex BEFORE \input{packages}.
\ifdraftmode
  \usepackage{lineno}\linenumbers   % line numbers for "see line 214" review comments (prose only; align/equation bodies aren't numbered without extra setup)
  \usepackage{showkeys}             % print \label / \ref keys in the margin
  \overfullrule=5pt                 % black bar in the margin marks every overfull line
\fi
.
\ifdraftmode
  \usepackage{lineno}\linenumbers   % line numbers for "see line 214" review comments (prose only; align/equation bodies aren't numbered without extra setup)
  \usepackage{showkeys}             % print \label / \ref keys in the margin
  \overfullrule=5pt                 % black bar in the margin marks every overfull line
\fi
.
\ifdraftmode
  \usepackage{lineno}\linenumbers   % line numbers for "see line 214" review comments (prose only; align/equation bodies aren't numbered without extra setup)
  \usepackage{showkeys}             % print \label / \ref keys in the margin
  \overfullrule=5pt                 % black bar in the margin marks every overfull line
\fi
.
\ifdraftmode
  \usepackage{lineno}\linenumbers   % line numbers for "see line 214" review comments (prose only; align/equation bodies aren't numbered without extra setup)
  \usepackage{showkeys}             % print \label / \ref keys in the margin
  \overfullrule=5pt                 % black bar in the margin marks every overfull line
\fi
